\documentclass[12pt]{article}

\usepackage{float}
\usepackage{graphicx}			% Use this package to include images
\usepackage{amsmath}			% A library of many standard math expressions
\usepackage[margin=1in]{geometry}% Sets 1in margins. 
\usepackage{fancyhdr}			% Creates headers and footers
\usepackage{enumerate}          %These two package give custom labels to a list
\usepackage[shortlabels]{enumitem}
\usepackage{matlab-prettifier}
\usepackage{tikz}
\usepackage[utf8]{inputenc}
\usepackage{pgfplots} 
\usepackage{pgfgantt}
\usepackage{pdflscape}
\usepackage{array,multirow}
\usepackage{tabularx,booktabs}
\newcolumntype{Y}{>{\centering\arraybackslash}X}
\usetikzlibrary{shapes.geometric}
\usetikzlibrary{arrows.meta,arrows}
\pgfplotsset{compat=newest} 
\pgfplotsset{plot coordinates/math parser=false}
\newlength\fwidth
\newlength\fheight
\usepackage{float,graphicx}
\pagestyle{fancy}
\fancyhead[l]{}
\fancyhead[c]{Experiment 9}
\fancyhead[r]{}
\fancyfoot[c]{\thepage}
\setlength{\headheight}{20pt}
\begin{document}
\author{}
\date{}
\title{Thermistor}
\maketitle
\section*{Aim}
To understand the working principle behind the thermistor.
\section*{Analysis}
The Steinhart-Hart Equation is used to convert resistance reading of a thermistor to thermodynamic temperature.
\[\frac{1}{T} = A + B\ln R + C(\ln R)^3\]
The formula was derived using fundamental principles of solid state physics. NTCs are polycrystalline
semiconductor devices; the resistance should follow the classical temperature activated form.
\[R = R_0 e^{B\big(\frac{1}{T}-\frac{1}{T_0}\big)}\]
which on solving gives
\[\frac{1}{T} = \frac{1}{T_0} + \frac{1}{B}\bigg( \ln \frac{R}{R_0}\bigg)\]
which dissolves into a linear relation
\[\frac{1}{T} = a_0 + a_1 \ln R\]
This is a good enough relation for small changes in resistance.
For larger deviations, $B$ itself becomes temperature dependent and we apply series expansion to get more accurate results.
\[\frac{1}{T} = \sum_{i =0 }^{n} a_i \ln ^i R\]
We find that $a_2$ is statistically insignificant and there is no loss in precision by terminating the expansion after $n =3 $.
This gives us back the Steinhart-Hart Equation.
We use both linear and cubic fits to approximate the relationship between $\ln R \text{ and } \frac{1}{T}$.
\vspace*{7pt}

\noindent \textbf{Safety precautions :}
The induction stove must not go beyond 800 watts.
The pipes carring steam must not be bent.
Use gloves to touch hot equipment.

\pagebreak

The data is different for heating and cooling the materials and follows a hysteresis.
Hence we perform the regression separately.
\begin{figure}[H]
    \centering
    \setlength\fwidth{0.7\textwidth}
    \setlength\fheight{0.5\textwidth}
    % This file was created by matlab2tikz.
%
\definecolor{mycolor1}{rgb}{0.00000,0.44700,0.74100}%
\definecolor{mycolor2}{rgb}{0.85000,0.32500,0.09800}%
\definecolor{mycolor3}{rgb}{0.12941,0.12941,0.12941}%
%
\begin{tikzpicture}

\begin{axis}[%
width=0.951\fwidth,
height=\fheight,
at={(0\fwidth,0\fheight)},
scale only axis,
xmin=-0.5,
xmax=1.35,
xlabel style={font=\color{mycolor3}},
xlabel={ln(R)},
ymin=0.0026,
ymax=0.0038,
ylabel style={font=\color{mycolor3}},
ylabel={1/T},
axis background/.style={fill=white},
title style={font=\bfseries\color{mycolor3}},
title={Heating},
xmajorgrids,
ymajorgrids,
legend style={at={(0.03,0.97)}, anchor=north west, legend cell align=left, align=left}
]
\addplot[only marks, mark=o, mark options={}, mark size=1.5000pt, draw=mycolor1] table[row sep=crcr]{%
x	y\\
1.03342375548695	0.00341296928327645\\
0.905795880367869	0.00337837837837838\\
0.837588438235511	0.00335570469798658\\
0.740362689494244	0.00332225913621262\\
0.68304703823885	0.0033003300330033\\
0.553883026643874	0.00324675324675325\\
0.41161970596323	0.00319488817891374\\
0.311753861055754	0.00314465408805031\\
0.225309281725863	0.00309597523219814\\
0.139879086401236	0.00304878048780488\\
0.0644579892269185	0.003003003003003\\
-0.0222763947111523	0.0029585798816568\\
-0.102372908709559	0.00291545189504373\\
-0.136677139879544	0.0028735632183908\\
-0.158015195409886	0.0028328611898017\\
-0.169411331314856	0.00280898876404494\\
};
\addlegendentry{Data}

\addplot [color=mycolor2]
  table[row sep=crcr]{%
-1	0.002508662\\
-0.996996996996997	0.00250903819257225\\
-0.993993993993994	0.00250942300523184\\
-0.990990990990991	0.00250981641201464\\
-0.987987987987988	0.00251021838695655\\
-0.984984984984985	0.00251062890409345\\
-0.981981981981982	0.0025110479374612\\
-0.978978978978979	0.0025114754610957\\
-0.975975975975976	0.00251191144903283\\
-0.972972972972973	0.00251235587530847\\
-0.96996996996997	0.0025128087139585\\
-0.966966966966967	0.0025132699390188\\
-0.963963963963964	0.00251373952452526\\
-0.960960960960961	0.00251421744451375\\
-0.957957957957958	0.00251470367302015\\
-0.954954954954955	0.00251519818408036\\
-0.951951951951952	0.00251570095173024\\
-0.948948948948949	0.00251621195000568\\
-0.945945945945946	0.00251673115294257\\
-0.942942942942943	0.00251725853457678\\
-0.93993993993994	0.0025177940689442\\
-0.936936936936937	0.0025183377300807\\
-0.933933933933934	0.00251888949202217\\
-0.930930930930931	0.0025194493288045\\
-0.927927927927928	0.00252001721446355\\
-0.924924924924925	0.00252059312303522\\
-0.921921921921922	0.00252117702855538\\
-0.918918918918919	0.00252176890505992\\
-0.915915915915916	0.00252236872658471\\
-0.912912912912913	0.00252297646716565\\
-0.90990990990991	0.0025235921008386\\
-0.906906906906907	0.00252421560163946\\
-0.903903903903904	0.0025248469436041\\
-0.900900900900901	0.00252548610076841\\
-0.897897897897898	0.00252613304716826\\
-0.894894894894895	0.00252678775683954\\
-0.891891891891892	0.00252745020381814\\
-0.888888888888889	0.00252812036213992\\
-0.885885885885886	0.00252879820584077\\
-0.882882882882883	0.00252948370895658\\
-0.87987987987988	0.00253017684552323\\
-0.876876876876877	0.00253087758957659\\
-0.873873873873874	0.00253158591515255\\
-0.870870870870871	0.00253230179628699\\
-0.867867867867868	0.0025330252070158\\
-0.864864864864865	0.00253375612137484\\
-0.861861861861862	0.00253449451340002\\
-0.858858858858859	0.00253524035712719\\
-0.855855855855856	0.00253599362659226\\
-0.852852852852853	0.0025367542958311\\
-0.84984984984985	0.00253752233887958\\
-0.846846846846847	0.0025382977297736\\
-0.843843843843844	0.00253908044254903\\
-0.840840840840841	0.00253987045124176\\
-0.837837837837838	0.00254066772988767\\
-0.834834834834835	0.00254147225252263\\
-0.831831831831832	0.00254228399318253\\
-0.828828828828829	0.00254310292590326\\
-0.825825825825826	0.00254392902472069\\
-0.822822822822823	0.0025447622636707\\
-0.81981981981982	0.00254560261678918\\
-0.816816816816817	0.002546450058112\\
-0.813813813813814	0.00254730456167506\\
-0.810810810810811	0.00254816610151422\\
-0.807807807807808	0.00254903465166538\\
-0.804804804804805	0.00254991018616441\\
-0.801801801801802	0.0025507926790472\\
-0.798798798798799	0.00255168210434962\\
-0.795795795795796	0.00255257843610756\\
-0.792792792792793	0.0025534816483569\\
-0.78978978978979	0.00255439171513352\\
-0.786786786786787	0.00255530861047331\\
-0.783783783783784	0.00255623230841214\\
-0.780780780780781	0.00255716278298589\\
-0.777777777777778	0.00255810000823045\\
-0.774774774774775	0.0025590439581817\\
-0.771771771771772	0.00255999460687552\\
-0.768768768768769	0.00256095192834779\\
-0.765765765765766	0.0025619158966344\\
-0.762762762762763	0.00256288648577122\\
-0.75975975975976	0.00256386366979414\\
-0.756756756756757	0.00256484742273903\\
-0.753753753753754	0.00256583771864178\\
-0.750750750750751	0.00256683453153827\\
-0.747747747747748	0.00256783783546439\\
-0.744744744744745	0.00256884760445601\\
-0.741741741741742	0.00256986381254902\\
-0.738738738738739	0.00257088643377929\\
-0.735735735735736	0.00257191544218271\\
-0.732732732732733	0.00257295081179516\\
-0.72972972972973	0.00257399251665252\\
-0.726726726726727	0.00257504053079067\\
-0.723723723723724	0.0025760948282455\\
-0.720720720720721	0.00257715538305288\\
-0.717717717717718	0.00257822216924871\\
-0.714714714714715	0.00257929516086885\\
-0.711711711711712	0.00258037433194919\\
-0.708708708708709	0.00258145965652561\\
-0.705705705705706	0.002582551108634\\
-0.702702702702703	0.00258364866231023\\
-0.6996996996997	0.00258475229159019\\
-0.696696696696697	0.00258586197050975\\
-0.693693693693694	0.00258697767310481\\
-0.690690690690691	0.00258809937341123\\
-0.687687687687688	0.00258922704546491\\
-0.684684684684685	0.00259036066330172\\
-0.681681681681682	0.00259150020095755\\
-0.678678678678679	0.00259264563246828\\
-0.675675675675676	0.00259379693186978\\
-0.672672672672673	0.00259495407319794\\
-0.66966966966967	0.00259611703048865\\
-0.666666666666667	0.00259728577777778\\
-0.663663663663664	0.00259846028910121\\
-0.660660660660661	0.00259964053849483\\
-0.657657657657658	0.00260082649999452\\
-0.654654654654655	0.00260201814763615\\
-0.651651651651652	0.00260321545545562\\
-0.648648648648649	0.0026044183974888\\
-0.645645645645646	0.00260562694777157\\
-0.642642642642643	0.00260684108033982\\
-0.63963963963964	0.00260806076922942\\
-0.636636636636637	0.00260928598847626\\
-0.633633633633634	0.00261051671211622\\
-0.630630630630631	0.00261175291418519\\
-0.627627627627628	0.00261299456871903\\
-0.624624624624625	0.00261424164975364\\
-0.621621621621622	0.0026154941313249\\
-0.618618618618619	0.00261675198746868\\
-0.615615615615616	0.00261801519222087\\
-0.612612612612613	0.00261928371961735\\
-0.60960960960961	0.00262055754369401\\
-0.606606606606607	0.00262183663848671\\
-0.603603603603604	0.00262312097803136\\
-0.600600600600601	0.00262441053636381\\
-0.597597597597598	0.00262570528751997\\
-0.594594594594595	0.0026270052055357\\
-0.591591591591592	0.0026283102644469\\
-0.588588588588589	0.00262962043828944\\
-0.585585585585586	0.0026309357010992\\
-0.582582582582583	0.00263225602691207\\
-0.57957957957958	0.00263358138976392\\
-0.576576576576577	0.00263491176369064\\
-0.573573573573574	0.00263624712272812\\
-0.570570570570571	0.00263758744091222\\
-0.567567567567568	0.00263893269227884\\
-0.564564564564564	0.00264028285086385\\
-0.561561561561562	0.00264163789070314\\
-0.558558558558559	0.00264299778583258\\
-0.555555555555556	0.00264436251028807\\
-0.552552552552553	0.00264573203810547\\
-0.549549549549549	0.00264710634332068\\
-0.546546546546547	0.00264848539996957\\
-0.543543543543544	0.00264986918208802\\
-0.540540540540541	0.00265125766371192\\
-0.537537537537538	0.00265265081887715\\
-0.534534534534534	0.00265404862161959\\
-0.531531531531532	0.00265545104597512\\
-0.528528528528529	0.00265685806597962\\
-0.525525525525526	0.00265826965566898\\
-0.522522522522523	0.00265968578907907\\
-0.519519519519519	0.00266110644024578\\
-0.516516516516517	0.00266253158320499\\
-0.513513513513513	0.00266396119199258\\
-0.510510510510511	0.00266539524064443\\
-0.507507507507508	0.00266683370319642\\
-0.504504504504504	0.00266827655368444\\
-0.501501501501502	0.00266972376614436\\
-0.498498498498498	0.00267117531461207\\
-0.495495495495496	0.00267263117312345\\
-0.492492492492492	0.00267409131571438\\
-0.489489489489489	0.00267555571642075\\
-0.486486486486487	0.00267702434927842\\
-0.483483483483483	0.0026784971883233\\
-0.480480480480481	0.00267997420759125\\
-0.477477477477477	0.00268145538111815\\
-0.474474474474474	0.0026829406829399\\
-0.471471471471472	0.00268443008709237\\
-0.468468468468468	0.00268592356761144\\
-0.465465465465465	0.00268742109853299\\
-0.462462462462462	0.00268892265389292\\
-0.459459459459459	0.00269042820772708\\
-0.456456456456456	0.00269193773407138\\
-0.453453453453453	0.00269345120696169\\
-0.45045045045045	0.00269496860043389\\
-0.447447447447447	0.00269648988852386\\
-0.444444444444444	0.00269801504526749\\
-0.441441441441441	0.00269954404470065\\
-0.438438438438438	0.00270107686085924\\
-0.435435435435435	0.00270261346777912\\
-0.432432432432432	0.00270415383949618\\
-0.429429429429429	0.0027056979500463\\
-0.426426426426426	0.00270724577346537\\
-0.423423423423423	0.00270879728378927\\
-0.42042042042042	0.00271035245505387\\
-0.417417417417417	0.00271191126129506\\
-0.414414414414414	0.00271347367654872\\
-0.411411411411411	0.00271503967485073\\
-0.408408408408408	0.00271660923023697\\
-0.405405405405405	0.00271818231674333\\
-0.402402402402402	0.00271975890840569\\
-0.399399399399399	0.00272133897925992\\
-0.396396396396396	0.00272292250334191\\
-0.393393393393393	0.00272450945468754\\
-0.39039039039039	0.0027260998073327\\
-0.387387387387387	0.00272769353531325\\
-0.384384384384384	0.0027292906126651\\
-0.381381381381381	0.00273089101342411\\
-0.378378378378378	0.00273249471162616\\
-0.375375375375375	0.00273410168130715\\
-0.372372372372372	0.00273571189650295\\
-0.369369369369369	0.00273732533124944\\
-0.366366366366366	0.0027389419595825\\
-0.363363363363363	0.00274056175553802\\
-0.36036036036036	0.00274218469315188\\
-0.357357357357357	0.00274381074645996\\
-0.354354354354354	0.00274543988949813\\
-0.351351351351351	0.00274707209630229\\
-0.348348348348348	0.00274870734090831\\
-0.345345345345345	0.00275034559735208\\
-0.342342342342342	0.00275198683966947\\
-0.339339339339339	0.00275363104189637\\
-0.336336336336336	0.00275527817806866\\
-0.333333333333333	0.00275692822222222\\
-0.33033033033033	0.00275858114839293\\
-0.327327327327327	0.00276023693061668\\
-0.324324324324324	0.00276189554292934\\
-0.321321321321321	0.0027635569593668\\
-0.318318318318318	0.00276522115396494\\
-0.315315315315315	0.00276688810075963\\
-0.312312312312312	0.00276855777378677\\
-0.309309309309309	0.00277023014708223\\
-0.306306306306306	0.0027719051946819\\
-0.303303303303303	0.00277358289062165\\
-0.3003003003003	0.00277526320893737\\
-0.297297297297297	0.00277694612366494\\
-0.294294294294294	0.00277863160884023\\
-0.291291291291291	0.00278031963849914\\
-0.288288288288288	0.00278201018667755\\
-0.285285285285285	0.00278370322741133\\
-0.282282282282282	0.00278539873473636\\
-0.279279279279279	0.00278709668268853\\
-0.276276276276276	0.00278879704530372\\
-0.273273273273273	0.00279049979661782\\
-0.27027027027027	0.00279220491066669\\
-0.267267267267267	0.00279391236148623\\
-0.264264264264264	0.00279562212311232\\
-0.261261261261261	0.00279733416958083\\
-0.258258258258258	0.00279904847492765\\
-0.255255255255255	0.00280076501318866\\
-0.252252252252252	0.00280248375839974\\
-0.249249249249249	0.00280420468459678\\
-0.246246246246246	0.00280592776581565\\
-0.243243243243243	0.00280765297609224\\
-0.24024024024024	0.00280938028946242\\
-0.237237237237237	0.00281110967996208\\
-0.234234234234234	0.00281284112162711\\
-0.231231231231231	0.00281457458849337\\
-0.228228228228228	0.00281631005459676\\
-0.225225225225225	0.00281804749397316\\
-0.222222222222222	0.00281978688065844\\
-0.219219219219219	0.00282152818868849\\
-0.216216216216216	0.00282327139209918\\
-0.213213213213213	0.00282501646492642\\
-0.21021021021021	0.00282676338120606\\
-0.207207207207207	0.002828512114974\\
-0.204204204204204	0.00283026264026611\\
-0.201201201201201	0.00283201493111828\\
-0.198198198198198	0.00283376896156639\\
-0.195195195195195	0.00283552470564632\\
-0.192192192192192	0.00283728213739395\\
-0.189189189189189	0.00283904123084516\\
-0.186186186186186	0.00284080196003584\\
-0.183183183183183	0.00284256429900187\\
-0.18018018018018	0.00284432822177912\\
-0.177177177177177	0.00284609370240348\\
-0.174174174174174	0.00284786071491084\\
-0.171171171171171	0.00284962923333706\\
-0.168168168168168	0.00285139923171804\\
-0.165165165165165	0.00285317068408966\\
-0.162162162162162	0.00285494356448779\\
-0.159159159159159	0.00285671784694832\\
-0.156156156156156	0.00285849350550713\\
-0.153153153153153	0.0028602705142001\\
-0.15015015015015	0.00286204884706312\\
-0.147147147147147	0.00286382847813206\\
-0.144144144144144	0.0028656093814428\\
-0.141141141141141	0.00286739153103123\\
-0.138138138138138	0.00286917490093324\\
-0.135135135135135	0.00287095946518469\\
-0.132132132132132	0.00287274519782147\\
-0.129129129129129	0.00287453207287947\\
-0.126126126126126	0.00287632006439456\\
-0.123123123123123	0.00287810914640263\\
-0.12012012012012	0.00287989929293956\\
-0.117117117117117	0.00288169047804123\\
-0.114114114114114	0.00288348267574351\\
-0.111111111111111	0.0028852758600823\\
-0.108108108108108	0.00288707000509348\\
-0.105105105105105	0.00288886508481292\\
-0.102102102102102	0.00289066107327651\\
-0.0990990990990991	0.00289245794452012\\
-0.0960960960960962	0.00289425567257965\\
-0.0930930930930931	0.00289605423149097\\
-0.0900900900900901	0.00289785359528996\\
-0.0870870870870871	0.0028996537380125\\
-0.0840840840840841	0.00290145463369449\\
-0.081081081081081	0.00290325625637178\\
-0.0780780780780781	0.00290505858008028\\
-0.075075075075075	0.00290686157885586\\
-0.0720720720720721	0.0029086652267344\\
-0.0690690690690691	0.00291046949775179\\
-0.066066066066066	0.0029122743659439\\
-0.0630630630630631	0.00291407980534662\\
-0.06006006006006	0.00291588578999582\\
-0.0570570570570571	0.0029176922939274\\
-0.0540540540540541	0.00291949929117723\\
-0.051051051051051	0.00292130675578119\\
-0.0480480480480481	0.00292311466177516\\
-0.045045045045045	0.00292492298319503\\
-0.0420420420420421	0.00292673169407668\\
-0.039039039039039	0.00292854076845598\\
-0.036036036036036	0.00293035018036883\\
-0.0330330330330331	0.0029321599038511\\
-0.03003003003003	0.00293396991293867\\
-0.027027027027027	0.00293578018166742\\
-0.024024024024024	0.00293759068407325\\
-0.021021021021021	0.00293940139419202\\
-0.0180180180180181	0.00294121228605962\\
-0.015015015015015	0.00294302333371193\\
-0.012012012012012	0.00294483451118483\\
-0.00900900900900903	0.00294664579251421\\
-0.00600600600600598	0.00294845715173594\\
-0.00300300300300305	0.00295026856288591\\
0	0.00295208\\
0.00300300300300305	0.00295389143711409\\
0.00600600600600609	0.00295570284826406\\
0.00900900900900892	0.00295751420748579\\
0.012012012012012	0.00295932548881517\\
0.015015015015015	0.00296113666628807\\
0.0180180180180181	0.00296294771394038\\
0.0210210210210211	0.00296475860580798\\
0.0240240240240239	0.00296656931592675\\
0.027027027027027	0.00296837981833258\\
0.03003003003003	0.00297019008706133\\
0.0330330330330331	0.0029720000961489\\
0.0360360360360361	0.00297380981963117\\
0.0390390390390389	0.00297561923154402\\
0.042042042042042	0.00297742830592332\\
0.045045045045045	0.00297923701680497\\
0.0480480480480481	0.00298104533822484\\
0.0510510510510511	0.00298285324421882\\
0.0540540540540539	0.00298466070882277\\
0.057057057057057	0.0029864677060726\\
0.06006006006006	0.00298827421000418\\
0.0630630630630631	0.00299008019465338\\
0.0660660660660661	0.0029918856340561\\
0.0690690690690692	0.00299369050224821\\
0.072072072072072	0.0029954947732656\\
0.075075075075075	0.00299729842114414\\
0.0780780780780781	0.00299910141991972\\
0.0810810810810811	0.00300090374362822\\
0.0840840840840842	0.00300270536630552\\
0.087087087087087	0.0030045062619875\\
0.0900900900900901	0.00300630640471004\\
0.0930930930930931	0.00300810576850903\\
0.0960960960960962	0.00300990432742035\\
0.0990990990990992	0.00301170205547988\\
0.102102102102102	0.00301349892672349\\
0.105105105105105	0.00301529491518708\\
0.108108108108108	0.00301708999490652\\
0.111111111111111	0.0030188841399177\\
0.114114114114114	0.00302067732425649\\
0.117117117117117	0.00302246952195877\\
0.12012012012012	0.00302426070706044\\
0.123123123123123	0.00302605085359737\\
0.126126126126126	0.00302783993560544\\
0.129129129129129	0.00302962792712053\\
0.132132132132132	0.00303141480217853\\
0.135135135135135	0.00303320053481531\\
0.138138138138138	0.00303498509906676\\
0.141141141141141	0.00303676846896877\\
0.144144144144144	0.0030385506185572\\
0.147147147147147	0.00304033152186794\\
0.15015015015015	0.00304211115293688\\
0.153153153153153	0.0030438894857999\\
0.156156156156156	0.00304566649449287\\
0.159159159159159	0.00304744215305168\\
0.162162162162162	0.00304921643551221\\
0.165165165165165	0.00305098931591034\\
0.168168168168168	0.00305276076828196\\
0.171171171171171	0.00305453076666294\\
0.174174174174174	0.00305629928508916\\
0.177177177177177	0.00305806629759652\\
0.18018018018018	0.00305983177822088\\
0.183183183183183	0.00306159570099813\\
0.186186186186186	0.00306335803996416\\
0.189189189189189	0.00306511876915484\\
0.192192192192192	0.00306687786260605\\
0.195195195195195	0.00306863529435368\\
0.198198198198198	0.00307039103843361\\
0.201201201201201	0.00307214506888172\\
0.204204204204204	0.00307389735973389\\
0.207207207207207	0.00307564788502601\\
0.21021021021021	0.00307739661879394\\
0.213213213213213	0.00307914353507359\\
0.216216216216216	0.00308088860790082\\
0.219219219219219	0.00308263181131151\\
0.222222222222222	0.00308437311934156\\
0.225225225225225	0.00308611250602684\\
0.228228228228228	0.00308784994540324\\
0.231231231231231	0.00308958541150663\\
0.234234234234234	0.00309131887837289\\
0.237237237237237	0.00309305032003792\\
0.24024024024024	0.00309477971053758\\
0.243243243243243	0.00309650702390776\\
0.246246246246246	0.00309823223418435\\
0.249249249249249	0.00309995531540322\\
0.252252252252252	0.00310167624160026\\
0.255255255255255	0.00310339498681134\\
0.258258258258258	0.00310511152507235\\
0.261261261261261	0.00310682583041917\\
0.264264264264264	0.00310853787688768\\
0.267267267267267	0.00311024763851377\\
0.27027027027027	0.00311195508933331\\
0.273273273273273	0.00311366020338218\\
0.276276276276276	0.00311536295469628\\
0.279279279279279	0.00311706331731147\\
0.282282282282282	0.00311876126526364\\
0.285285285285285	0.00312045677258868\\
0.288288288288288	0.00312214981332245\\
0.291291291291291	0.00312384036150086\\
0.294294294294294	0.00312552839115977\\
0.297297297297297	0.00312721387633506\\
0.3003003003003	0.00312889679106263\\
0.303303303303303	0.00313057710937835\\
0.306306306306306	0.0031322548053181\\
0.309309309309309	0.00313392985291777\\
0.312312312312312	0.00313560222621323\\
0.315315315315315	0.00313727189924037\\
0.318318318318318	0.00313893884603506\\
0.321321321321321	0.0031406030406332\\
0.324324324324324	0.00314226445707066\\
0.327327327327327	0.00314392306938332\\
0.33033033033033	0.00314557885160707\\
0.333333333333333	0.00314723177777778\\
0.336336336336336	0.00314888182193134\\
0.339339339339339	0.00315052895810363\\
0.342342342342342	0.00315217316033053\\
0.345345345345345	0.00315381440264792\\
0.348348348348348	0.00315545265909169\\
0.351351351351351	0.00315708790369771\\
0.354354354354354	0.00315872011050187\\
0.357357357357357	0.00316034925354004\\
0.36036036036036	0.00316197530684812\\
0.363363363363363	0.00316359824446198\\
0.366366366366366	0.0031652180404175\\
0.369369369369369	0.00316683466875056\\
0.372372372372372	0.00316844810349705\\
0.375375375375375	0.00317005831869285\\
0.378378378378378	0.00317166528837384\\
0.381381381381381	0.0031732689865759\\
0.384384384384384	0.0031748693873349\\
0.387387387387387	0.00317646646468675\\
0.39039039039039	0.0031780601926673\\
0.393393393393393	0.00317965054531246\\
0.396396396396396	0.00318123749665809\\
0.399399399399399	0.00318282102074008\\
0.402402402402402	0.00318440109159431\\
0.405405405405405	0.00318597768325667\\
0.408408408408408	0.00318755076976303\\
0.411411411411411	0.00318912032514927\\
0.414414414414414	0.00319068632345128\\
0.417417417417417	0.00319224873870494\\
0.42042042042042	0.00319380754494613\\
0.423423423423423	0.00319536271621073\\
0.426426426426426	0.00319691422653463\\
0.429429429429429	0.0031984620499537\\
0.432432432432432	0.00320000616050382\\
0.435435435435436	0.00320154653222088\\
0.438438438438439	0.00320308313914076\\
0.441441441441441	0.00320461595529935\\
0.444444444444444	0.00320614495473251\\
0.447447447447447	0.00320767011147614\\
0.450450450450451	0.00320919139956611\\
0.453453453453454	0.00321070879303831\\
0.456456456456456	0.00321222226592862\\
0.459459459459459	0.00321373179227292\\
0.462462462462462	0.00321523734610708\\
0.465465465465466	0.00321673890146701\\
0.468468468468469	0.00321823643238856\\
0.471471471471471	0.00321972991290763\\
0.474474474474474	0.0032212193170601\\
0.477477477477477	0.00322270461888185\\
0.480480480480481	0.00322418579240876\\
0.483483483483484	0.0032256628116767\\
0.486486486486486	0.00322713565072158\\
0.489489489489489	0.00322860428357925\\
0.492492492492492	0.00323006868428562\\
0.495495495495496	0.00323152882687655\\
0.498498498498499	0.00323298468538793\\
0.501501501501501	0.00323443623385564\\
0.504504504504504	0.00323588344631556\\
0.507507507507508	0.00323732629680358\\
0.510510510510511	0.00323876475935557\\
0.513513513513514	0.00324019880800742\\
0.516516516516516	0.00324162841679501\\
0.519519519519519	0.00324305355975422\\
0.522522522522523	0.00324447421092093\\
0.525525525525526	0.00324589034433102\\
0.528528528528529	0.00324730193402038\\
0.531531531531531	0.00324870895402488\\
0.534534534534534	0.00325011137838041\\
0.537537537537538	0.00325150918112285\\
0.540540540540541	0.00325290233628808\\
0.543543543543544	0.00325429081791198\\
0.546546546546546	0.00325567460003043\\
0.549549549549549	0.00325705365667932\\
0.552552552552553	0.00325842796189453\\
0.555555555555556	0.00325979748971193\\
0.558558558558559	0.00326116221416742\\
0.561561561561561	0.00326252210929686\\
0.564564564564564	0.00326387714913615\\
0.567567567567568	0.00326522730772116\\
0.570570570570571	0.00326657255908778\\
0.573573573573574	0.00326791287727188\\
0.576576576576577	0.00326924823630936\\
0.57957957957958	0.00327057861023608\\
0.582582582582583	0.00327190397308793\\
0.585585585585586	0.0032732242989008\\
0.588588588588589	0.00327453956171056\\
0.591591591591592	0.0032758497355531\\
0.594594594594595	0.0032771547944643\\
0.597597597597598	0.00327845471248003\\
0.600600600600601	0.00327974946363619\\
0.603603603603604	0.00328103902196865\\
0.606606606606607	0.00328232336151329\\
0.60960960960961	0.00328360245630599\\
0.612612612612613	0.00328487628038265\\
0.615615615615616	0.00328614480777913\\
0.618618618618619	0.00328740801253132\\
0.621621621621622	0.0032886658686751\\
0.624624624624625	0.00328991835024636\\
0.627627627627628	0.00329116543128097\\
0.630630630630631	0.00329240708581481\\
0.633633633633634	0.00329364328788378\\
0.636636636636637	0.00329487401152374\\
0.63963963963964	0.00329609923077058\\
0.642642642642643	0.00329731891966018\\
0.645645645645646	0.00329853305222843\\
0.648648648648649	0.0032997416025112\\
0.651651651651652	0.00330094454454438\\
0.654654654654655	0.00330214185236385\\
0.657657657657658	0.00330333350000548\\
0.660660660660661	0.00330451946150517\\
0.663663663663664	0.00330569971089879\\
0.666666666666667	0.00330687422222222\\
0.66966966966967	0.00330804296951135\\
0.672672672672673	0.00330920592680206\\
0.675675675675676	0.00331036306813022\\
0.678678678678679	0.00331151436753172\\
0.681681681681682	0.00331265979904245\\
0.684684684684685	0.00331379933669828\\
0.687687687687688	0.00331493295453509\\
0.690690690690691	0.00331606062658877\\
0.693693693693694	0.00331718232689519\\
0.696696696696697	0.00331829802949025\\
0.6996996996997	0.00331940770840981\\
0.702702702702703	0.00332051133768977\\
0.705705705705706	0.003321608891366\\
0.708708708708709	0.00332270034347439\\
0.711711711711712	0.00332378566805081\\
0.714714714714715	0.00332486483913115\\
0.717717717717718	0.00332593783075129\\
0.720720720720721	0.00332700461694712\\
0.723723723723724	0.0033280651717545\\
0.726726726726727	0.00332911946920933\\
0.72972972972973	0.00333016748334748\\
0.732732732732733	0.00333120918820484\\
0.735735735735736	0.00333224455781729\\
0.738738738738739	0.00333327356622071\\
0.741741741741742	0.00333429618745098\\
0.744744744744745	0.00333531239554399\\
0.747747747747748	0.00333632216453561\\
0.750750750750751	0.00333732546846173\\
0.753753753753754	0.00333832228135822\\
0.756756756756757	0.00333931257726097\\
0.75975975975976	0.00334029633020587\\
0.762762762762763	0.00334127351422878\\
0.765765765765766	0.0033422441033656\\
0.768768768768769	0.00334320807165221\\
0.771771771771772	0.00334416539312448\\
0.774774774774775	0.0033451160418183\\
0.777777777777778	0.00334605999176955\\
0.780780780780781	0.00334699721701411\\
0.783783783783784	0.00334792769158786\\
0.786786786786787	0.00334885138952669\\
0.78978978978979	0.00334976828486648\\
0.792792792792793	0.0033506783516431\\
0.795795795795796	0.00335158156389244\\
0.798798798798799	0.00335247789565038\\
0.801801801801802	0.0033533673209528\\
0.804804804804805	0.00335424981383559\\
0.807807807807808	0.00335512534833462\\
0.810810810810811	0.00335599389848578\\
0.813813813813814	0.00335685543832494\\
0.816816816816817	0.003357709941888\\
0.81981981981982	0.00335855738321082\\
0.822822822822823	0.0033593977363293\\
0.825825825825826	0.00336023097527931\\
0.828828828828829	0.00336105707409674\\
0.831831831831832	0.00336187600681747\\
0.834834834834835	0.00336268774747737\\
0.837837837837838	0.00336349227011233\\
0.840840840840841	0.00336428954875824\\
0.843843843843844	0.00336507955745097\\
0.846846846846847	0.0033658622702264\\
0.84984984984985	0.00336663766112042\\
0.852852852852853	0.0033674057041689\\
0.855855855855856	0.00336816637340774\\
0.858858858858859	0.00336891964287281\\
0.861861861861862	0.00336966548659998\\
0.864864864864865	0.00337040387862516\\
0.867867867867868	0.0033711347929842\\
0.870870870870871	0.00337185820371301\\
0.873873873873874	0.00337257408484745\\
0.876876876876877	0.00337328241042341\\
0.87987987987988	0.00337398315447677\\
0.882882882882883	0.00337467629104342\\
0.885885885885886	0.00337536179415923\\
0.888888888888889	0.00337603963786008\\
0.891891891891892	0.00337670979618187\\
0.894894894894895	0.00337737224316046\\
0.897897897897898	0.00337802695283174\\
0.900900900900901	0.00337867389923159\\
0.903903903903904	0.0033793130563959\\
0.906906906906907	0.00337994439836054\\
0.90990990990991	0.0033805678991614\\
0.912912912912913	0.00338118353283435\\
0.915915915915916	0.00338179127341529\\
0.918918918918919	0.00338239109494008\\
0.921921921921922	0.00338298297144462\\
0.924924924924925	0.00338356687696478\\
0.927927927927928	0.00338414278553645\\
0.930930930930931	0.0033847106711955\\
0.933933933933934	0.00338527050797783\\
0.936936936936937	0.0033858222699193\\
0.93993993993994	0.0033863659310558\\
0.942942942942943	0.00338690146542322\\
0.945945945945946	0.00338742884705743\\
0.948948948948949	0.00338794804999432\\
0.951951951951952	0.00338845904826976\\
0.954954954954955	0.00338896181591965\\
0.957957957957958	0.00338945632697985\\
0.960960960960961	0.00338994255548626\\
0.963963963963964	0.00339042047547474\\
0.966966966966967	0.0033908900609812\\
0.96996996996997	0.0033913512860415\\
0.972972972972973	0.00339180412469153\\
0.975975975975976	0.00339224855096717\\
0.978978978978979	0.0033926845389043\\
0.981981981981982	0.0033931120625388\\
0.984984984984985	0.00339353109590656\\
0.987987987987988	0.00339394161304345\\
0.990990990990991	0.00339434358798536\\
0.993993993993994	0.00339473699476816\\
0.996996996996997	0.00339512180742775\\
1	0.003395498\\
1.003003003003	0.00339586554652079\\
1.00600600600601	0.00339622442102601\\
1.00900900900901	0.00339657459755153\\
1.01201201201201	0.00339691605013324\\
1.01501501501502	0.00339724875280702\\
1.01801801801802	0.00339757267960875\\
1.02102102102102	0.00339788780457432\\
1.02402402402402	0.00339819410173959\\
1.02702702702703	0.00339849154514047\\
1.03003003003003	0.00339878010881281\\
1.03303303303303	0.00339905976679252\\
1.03603603603604	0.00339933049311547\\
1.03903903903904	0.00339959226181754\\
1.04204204204204	0.00339984504693461\\
1.04504504504505	0.00340008882250256\\
1.04804804804805	0.00340032356255728\\
1.05105105105105	0.00340054924113465\\
1.05405405405405	0.00340076583227055\\
1.05705705705706	0.00340097331000085\\
1.06006006006006	0.00340117164836145\\
1.06306306306306	0.00340136082138823\\
1.06606606606607	0.00340154080311705\\
1.06906906906907	0.00340171156758382\\
1.07207207207207	0.0034018730888244\\
1.07507507507508	0.00340202534087468\\
1.07807807807808	0.00340216829777054\\
1.08108108108108	0.00340230193354787\\
1.08408408408408	0.00340242622224253\\
1.08708708708709	0.00340254113789043\\
1.09009009009009	0.00340264665452743\\
1.09309309309309	0.00340274274618942\\
1.0960960960961	0.00340282938691228\\
1.0990990990991	0.00340290655073189\\
1.1021021021021	0.00340297421168413\\
1.10510510510511	0.00340303234380489\\
1.10810810810811	0.00340308092113004\\
1.11111111111111	0.00340311991769547\\
1.11411411411411	0.00340314930753706\\
1.11711711711712	0.0034031690646907\\
1.12012012012012	0.00340317916319225\\
1.12312312312312	0.00340317957707761\\
1.12612612612613	0.00340317028038265\\
1.12912912912913	0.00340315124714326\\
1.13213213213213	0.00340312245139531\\
1.13513513513514	0.0034030838671747\\
1.13813813813814	0.0034030354685173\\
1.14114114114114	0.00340297722945899\\
1.14414414414414	0.00340290912403565\\
1.14714714714715	0.00340283112628317\\
1.15015015015015	0.00340274321023743\\
1.15315315315315	0.0034026453499343\\
1.15615615615616	0.00340253751940968\\
1.15915915915916	0.00340241969269944\\
1.16216216216216	0.00340229184383946\\
1.16516516516517	0.00340215394686562\\
1.16816816816817	0.00340200597581381\\
1.17117117117117	0.00340184790471991\\
1.17417417417417	0.0034016797076198\\
1.17717717717718	0.00340150135854936\\
1.18018018018018	0.00340131283154447\\
1.18318318318318	0.00340111410064102\\
1.18618618618619	0.00340090513987488\\
1.18918918918919	0.00340068592328194\\
1.19219219219219	0.00340045642489807\\
1.1951951951952	0.00340021661875917\\
1.1981981981982	0.00339996647890111\\
1.2012012012012	0.00339970597935977\\
1.2042042042042	0.00339943509417103\\
1.20720720720721	0.00339915379737078\\
1.21021021021021	0.0033988620629949\\
1.21321321321321	0.00339855986507927\\
1.21621621621622	0.00339824717765976\\
1.21921921921922	0.00339792397477227\\
1.22222222222222	0.00339759023045267\\
1.22522522522523	0.00339724591873685\\
1.22822822822823	0.00339689101366069\\
1.23123123123123	0.00339652548926006\\
1.23423423423423	0.00339614931957085\\
1.23723723723724	0.00339576247862894\\
1.24024024024024	0.00339536494047022\\
1.24324324324324	0.00339495667913055\\
1.24624624624625	0.00339453766864584\\
1.24924924924925	0.00339410788305195\\
1.25225225225225	0.00339366729638477\\
1.25525525525526	0.00339321588268018\\
1.25825825825826	0.00339275361597406\\
1.26126126126126	0.0033922804703023\\
1.26426426426426	0.00339179641970077\\
1.26726726726727	0.00339130143820535\\
1.27027027027027	0.00339079549985193\\
1.27327327327327	0.0033902785786764\\
1.27627627627628	0.00338975064871462\\
1.27927927927928	0.00338921168400248\\
1.28228228228228	0.00338866165857587\\
1.28528528528529	0.00338810054647067\\
1.28828828828829	0.00338752832172275\\
1.29129129129129	0.00338694495836799\\
1.29429429429429	0.00338635043044229\\
1.2972972972973	0.00338574471198152\\
1.3003003003003	0.00338512777702156\\
1.3033033033033	0.0033844995995983\\
1.30630630630631	0.00338386015374761\\
1.30930930930931	0.00338320941350538\\
1.31231231231231	0.00338254735290749\\
1.31531531531532	0.00338187394598982\\
1.31831831831832	0.00338118916678825\\
1.32132132132132	0.00338049298933866\\
1.32432432432432	0.00337978538767694\\
1.32732732732733	0.00337906633583896\\
1.33033033033033	0.00337833580786062\\
1.33333333333333	0.00337759377777778\\
1.33633633633634	0.00337684021962633\\
1.33933933933934	0.00337607510744216\\
1.34234234234234	0.00337529841526113\\
1.34534534534535	0.00337451011711915\\
1.34834834834835	0.00337371018705208\\
1.35135135135135	0.00337289859909581\\
1.35435435435435	0.00337207532728622\\
1.35735735735736	0.00337124034565919\\
1.36036036036036	0.0033703936282506\\
1.36336336336336	0.00336953514909634\\
1.36636636636637	0.00336866488223229\\
1.36936936936937	0.00336778280169432\\
1.37237237237237	0.00336688888151832\\
1.37537537537538	0.00336598309574017\\
1.37837837837838	0.00336506541839575\\
1.38138138138138	0.00336413582352095\\
1.38438438438438	0.00336319428515164\\
1.38738738738739	0.00336224077732371\\
1.39039039039039	0.00336127527407303\\
1.39339339339339	0.0033602977494355\\
1.3963963963964	0.00335930817744699\\
1.3993993993994	0.00335830653214338\\
1.4024024024024	0.00335729278756055\\
1.40540540540541	0.00335626691773439\\
1.40840840840841	0.00335522889670078\\
1.41141141141141	0.00335417869849559\\
1.41441441441441	0.00335311629715471\\
1.41741741741742	0.00335204166671403\\
1.42042042042042	0.00335095478120942\\
1.42342342342342	0.00334985561467677\\
1.42642642642643	0.00334874414115195\\
1.42942942942943	0.00334762033467084\\
1.43243243243243	0.00334648416926934\\
1.43543543543544	0.00334533561898332\\
1.43843843843844	0.00334417465784866\\
1.44144144144144	0.00334300125990125\\
1.44444444444444	0.00334181539917696\\
1.44744744744745	0.00334061704971167\\
1.45045045045045	0.00333940618554128\\
1.45345345345345	0.00333818278070165\\
1.45645645645646	0.00333694680922868\\
1.45945945945946	0.00333569824515823\\
1.46246246246246	0.00333443706252621\\
1.46546546546547	0.00333316323536847\\
1.46846846846847	0.00333187673772092\\
1.47147147147147	0.00333057754361943\\
1.47447447447447	0.00332926562709987\\
1.47747747747748	0.00332794096219814\\
1.48048048048048	0.00332660352295011\\
1.48348348348348	0.00332525328339166\\
1.48648648648649	0.00332389021755868\\
1.48948948948949	0.00332251429948705\\
1.49249249249249	0.00332112550321265\\
1.4954954954955	0.00331972380277136\\
1.4984984984985	0.00331830917219906\\
1.5015015015015	0.00331688158553164\\
1.5045045045045	0.00331544101680497\\
1.50750750750751	0.00331398744005494\\
1.51051051051051	0.00331252082931743\\
1.51351351351351	0.00331104115862831\\
1.51651651651652	0.00330954840202348\\
1.51951951951952	0.00330804253353881\\
1.52252252252252	0.00330652352721019\\
1.52552552552553	0.00330499135707349\\
1.52852852852853	0.0033034459971646\\
1.53153153153153	0.0033018874215194\\
1.53453453453453	0.00330031560417377\\
1.53753753753754	0.00329873051916359\\
1.54054054054054	0.00329713214052475\\
1.54354354354354	0.00329552044229312\\
1.54654654654655	0.00329389539850458\\
1.54954954954955	0.00329225698319503\\
1.55255255255255	0.00329060517040033\\
1.55555555555556	0.00328893993415638\\
1.55855855855856	0.00328726124849905\\
1.56156156156156	0.00328556908746422\\
1.56456456456456	0.00328386342508778\\
1.56756756756757	0.0032821442354056\\
1.57057057057057	0.00328041149245358\\
1.57357357357357	0.00327866517026758\\
1.57657657657658	0.0032769052428835\\
1.57957957957958	0.00327513168433721\\
1.58258258258258	0.00327334446866459\\
1.58558558558559	0.00327154356990153\\
1.58858858858859	0.00326972896208391\\
1.59159159159159	0.00326790061924761\\
1.59459459459459	0.0032660585154285\\
1.5975975975976	0.00326420262466248\\
1.6006006006006	0.00326233292098543\\
1.6036036036036	0.00326044937843322\\
1.60660660660661	0.00325855197104173\\
1.60960960960961	0.00325664067284686\\
1.61261261261261	0.00325471545788447\\
1.61561561561562	0.00325277630019046\\
1.61861861861862	0.0032508231738007\\
1.62162162162162	0.00324885605275107\\
1.62462462462462	0.00324687491107746\\
1.62762762762763	0.00324487972281575\\
1.63063063063063	0.00324287046200181\\
1.63363363363363	0.00324084710267154\\
1.63663663663664	0.00323880961886081\\
1.63963963963964	0.0032367579846055\\
1.64264264264264	0.00323469217394149\\
1.64564564564565	0.00323261216090468\\
1.64864864864865	0.00323051791953093\\
1.65165165165165	0.00322840942385613\\
1.65465465465465	0.00322628664791616\\
1.65765765765766	0.0032241495657469\\
1.66066066066066	0.00322199815138424\\
1.66366366366366	0.00321983237886405\\
1.66666666666667	0.00321765222222222\\
1.66966966966967	0.00321545765549463\\
1.67267267267267	0.00321324865271716\\
1.67567567567568	0.00321102518792569\\
1.67867867867868	0.0032087872351561\\
1.68168168168168	0.00320653476844428\\
1.68468468468468	0.00320426776182611\\
1.68768768768769	0.00320198618933746\\
1.69069069069069	0.00319969002501422\\
1.69369369369369	0.00319737924289227\\
1.6966966966967	0.0031950538170075\\
1.6996996996997	0.00319271372139577\\
1.7027027027027	0.00319035893009299\\
1.70570570570571	0.00318798941713502\\
1.70870870870871	0.00318560515655774\\
1.71171171171171	0.00318320612239705\\
1.71471471471471	0.00318079228868882\\
1.71771771771772	0.00317836362946893\\
1.72072072072072	0.00317592011877327\\
1.72372372372372	0.00317346173063771\\
1.72672672672673	0.00317098843909814\\
1.72972972972973	0.00316850021819043\\
1.73273273273273	0.00316599704195048\\
1.73573573573574	0.00316347888441416\\
1.73873873873874	0.00316094571961735\\
1.74174174174174	0.00315839752159594\\
1.74474474474474	0.0031558342643858\\
1.74774774774775	0.00315325592202283\\
1.75075075075075	0.00315066246854289\\
1.75375375375375	0.00314805387798187\\
1.75675675675676	0.00314543012437565\\
1.75975975975976	0.00314279118176012\\
1.76276276276276	0.00314013702417116\\
1.76576576576577	0.00313746762564464\\
1.76876876876877	0.00313478296021645\\
1.77177177177177	0.00313208300192247\\
1.77477477477477	0.00312936772479858\\
1.77777777777778	0.00312663710288066\\
1.78078078078078	0.0031238911102046\\
1.78378378378378	0.00312112972080627\\
1.78678678678679	0.00311835290872156\\
1.78978978978979	0.00311556064798635\\
1.79279279279279	0.00311275291263652\\
1.7957957957958	0.00310992967670796\\
1.7987987987988	0.00310709091423653\\
1.8018018018018	0.00310423659925813\\
1.8048048048048	0.00310136670580864\\
1.80780780780781	0.00309848120792394\\
1.81081081081081	0.0030955800796399\\
1.81381381381381	0.00309266329499242\\
1.81681681681682	0.00308973082801737\\
1.81981981981982	0.00308678265275063\\
1.82282282282282	0.00308381874322809\\
1.82582582582583	0.00308083907348563\\
1.82882882882883	0.00307784361755912\\
1.83183183183183	0.00307483234948446\\
1.83483483483483	0.00307180524329751\\
1.83783783783784	0.00306876227303417\\
1.84084084084084	0.00306570341273032\\
1.84384384384384	0.00306262863642183\\
1.84684684684685	0.00305953791814459\\
1.84984984984985	0.00305643123193447\\
1.85285285285285	0.00305330855182737\\
1.85585585585586	0.00305016985185916\\
1.85885885885886	0.00304701510606573\\
1.86186186186186	0.00304384428848295\\
1.86486486486486	0.0030406573731467\\
1.86786786786787	0.00303745433409288\\
1.87087087087087	0.00303423514535735\\
1.87387387387387	0.00303099978097601\\
1.87687687687688	0.00302774821498473\\
1.87987987987988	0.00302448042141939\\
1.88288288288288	0.00302119637431588\\
1.88588588588589	0.00301789604771007\\
1.88888888888889	0.00301457941563786\\
1.89189189189189	0.00301124645213511\\
1.89489489489489	0.00300789713123772\\
1.8978978978979	0.00300453142698156\\
1.9009009009009	0.00300114931340252\\
1.9039039039039	0.00299775076453647\\
1.90690690690691	0.0029943357544193\\
1.90990990990991	0.00299090425708689\\
1.91291291291291	0.00298745624657512\\
1.91591591591592	0.00298399169691987\\
1.91891891891892	0.00298051058215703\\
1.92192192192192	0.00297701287632247\\
1.92492492492492	0.00297349855345208\\
1.92792792792793	0.00296996758758174\\
1.93093093093093	0.00296641995274733\\
1.93393393393393	0.00296285562298472\\
1.93693693693694	0.00295927457232982\\
1.93993993993994	0.00295567677481848\\
1.94294294294294	0.00295206220448661\\
1.94594594594595	0.00294843083537007\\
1.94894894894895	0.00294478264150475\\
1.95195195195195	0.00294111759692653\\
1.95495495495495	0.00293743567567129\\
1.95795795795796	0.00293373685177491\\
1.96096096096096	0.00293002109927329\\
1.96396396396396	0.00292628839220228\\
1.96696696696697	0.00292253870459779\\
1.96996996996997	0.00291877201049568\\
1.97297297297297	0.00291498828393185\\
1.97597597597598	0.00291118749894217\\
1.97897897897898	0.00290736962956252\\
1.98198198198198	0.00290353464982879\\
1.98498498498498	0.00289968253377686\\
1.98798798798799	0.0028958132554426\\
1.99099099099099	0.00289192678886191\\
1.99399399399399	0.00288802310807066\\
1.996996996997	0.00288410218710473\\
2	0.002880164\\
};
\addlegendentry{Fit}

\addplot [color=black, dashed]
  table[row sep=crcr]{%
-1	0.002465551\\
-0.996996996996997	0.00246702222822823\\
-0.993993993993994	0.00246849345645646\\
-0.990990990990991	0.00246996468468468\\
-0.987987987987988	0.00247143591291291\\
-0.984984984984985	0.00247290714114114\\
-0.981981981981982	0.00247437836936937\\
-0.978978978978979	0.0024758495975976\\
-0.975975975975976	0.00247732082582583\\
-0.972972972972973	0.00247879205405405\\
-0.96996996996997	0.00248026328228228\\
-0.966966966966967	0.00248173451051051\\
-0.963963963963964	0.00248320573873874\\
-0.960960960960961	0.00248467696696697\\
-0.957957957957958	0.0024861481951952\\
-0.954954954954955	0.00248761942342342\\
-0.951951951951952	0.00248909065165165\\
-0.948948948948949	0.00249056187987988\\
-0.945945945945946	0.00249203310810811\\
-0.942942942942943	0.00249350433633634\\
-0.93993993993994	0.00249497556456456\\
-0.936936936936937	0.00249644679279279\\
-0.933933933933934	0.00249791802102102\\
-0.930930930930931	0.00249938924924925\\
-0.927927927927928	0.00250086047747748\\
-0.924924924924925	0.00250233170570571\\
-0.921921921921922	0.00250380293393393\\
-0.918918918918919	0.00250527416216216\\
-0.915915915915916	0.00250674539039039\\
-0.912912912912913	0.00250821661861862\\
-0.90990990990991	0.00250968784684685\\
-0.906906906906907	0.00251115907507508\\
-0.903903903903904	0.0025126303033033\\
-0.900900900900901	0.00251410153153153\\
-0.897897897897898	0.00251557275975976\\
-0.894894894894895	0.00251704398798799\\
-0.891891891891892	0.00251851521621622\\
-0.888888888888889	0.00251998644444444\\
-0.885885885885886	0.00252145767267267\\
-0.882882882882883	0.0025229289009009\\
-0.87987987987988	0.00252440012912913\\
-0.876876876876877	0.00252587135735736\\
-0.873873873873874	0.00252734258558559\\
-0.870870870870871	0.00252881381381381\\
-0.867867867867868	0.00253028504204204\\
-0.864864864864865	0.00253175627027027\\
-0.861861861861862	0.0025332274984985\\
-0.858858858858859	0.00253469872672673\\
-0.855855855855856	0.00253616995495495\\
-0.852852852852853	0.00253764118318318\\
-0.84984984984985	0.00253911241141141\\
-0.846846846846847	0.00254058363963964\\
-0.843843843843844	0.00254205486786787\\
-0.840840840840841	0.0025435260960961\\
-0.837837837837838	0.00254499732432432\\
-0.834834834834835	0.00254646855255255\\
-0.831831831831832	0.00254793978078078\\
-0.828828828828829	0.00254941100900901\\
-0.825825825825826	0.00255088223723724\\
-0.822822822822823	0.00255235346546547\\
-0.81981981981982	0.00255382469369369\\
-0.816816816816817	0.00255529592192192\\
-0.813813813813814	0.00255676715015015\\
-0.810810810810811	0.00255823837837838\\
-0.807807807807808	0.00255970960660661\\
-0.804804804804805	0.00256118083483483\\
-0.801801801801802	0.00256265206306306\\
-0.798798798798799	0.00256412329129129\\
-0.795795795795796	0.00256559451951952\\
-0.792792792792793	0.00256706574774775\\
-0.78978978978979	0.00256853697597598\\
-0.786786786786787	0.0025700082042042\\
-0.783783783783784	0.00257147943243243\\
-0.780780780780781	0.00257295066066066\\
-0.777777777777778	0.00257442188888889\\
-0.774774774774775	0.00257589311711712\\
-0.771771771771772	0.00257736434534535\\
-0.768768768768769	0.00257883557357357\\
-0.765765765765766	0.0025803068018018\\
-0.762762762762763	0.00258177803003003\\
-0.75975975975976	0.00258324925825826\\
-0.756756756756757	0.00258472048648649\\
-0.753753753753754	0.00258619171471471\\
-0.750750750750751	0.00258766294294294\\
-0.747747747747748	0.00258913417117117\\
-0.744744744744745	0.0025906053993994\\
-0.741741741741742	0.00259207662762763\\
-0.738738738738739	0.00259354785585586\\
-0.735735735735736	0.00259501908408408\\
-0.732732732732733	0.00259649031231231\\
-0.72972972972973	0.00259796154054054\\
-0.726726726726727	0.00259943276876877\\
-0.723723723723724	0.002600903996997\\
-0.720720720720721	0.00260237522522523\\
-0.717717717717718	0.00260384645345345\\
-0.714714714714715	0.00260531768168168\\
-0.711711711711712	0.00260678890990991\\
-0.708708708708709	0.00260826013813814\\
-0.705705705705706	0.00260973136636637\\
-0.702702702702703	0.00261120259459459\\
-0.6996996996997	0.00261267382282282\\
-0.696696696696697	0.00261414505105105\\
-0.693693693693694	0.00261561627927928\\
-0.690690690690691	0.00261708750750751\\
-0.687687687687688	0.00261855873573574\\
-0.684684684684685	0.00262002996396396\\
-0.681681681681682	0.00262150119219219\\
-0.678678678678679	0.00262297242042042\\
-0.675675675675676	0.00262444364864865\\
-0.672672672672673	0.00262591487687688\\
-0.66966966966967	0.0026273861051051\\
-0.666666666666667	0.00262885733333333\\
-0.663663663663664	0.00263032856156156\\
-0.660660660660661	0.00263179978978979\\
-0.657657657657658	0.00263327101801802\\
-0.654654654654655	0.00263474224624625\\
-0.651651651651652	0.00263621347447447\\
-0.648648648648649	0.0026376847027027\\
-0.645645645645646	0.00263915593093093\\
-0.642642642642643	0.00264062715915916\\
-0.63963963963964	0.00264209838738739\\
-0.636636636636637	0.00264356961561562\\
-0.633633633633634	0.00264504084384384\\
-0.630630630630631	0.00264651207207207\\
-0.627627627627628	0.0026479833003003\\
-0.624624624624625	0.00264945452852853\\
-0.621621621621622	0.00265092575675676\\
-0.618618618618619	0.00265239698498499\\
-0.615615615615616	0.00265386821321321\\
-0.612612612612613	0.00265533944144144\\
-0.60960960960961	0.00265681066966967\\
-0.606606606606607	0.0026582818978979\\
-0.603603603603604	0.00265975312612613\\
-0.600600600600601	0.00266122435435435\\
-0.597597597597598	0.00266269558258258\\
-0.594594594594595	0.00266416681081081\\
-0.591591591591592	0.00266563803903904\\
-0.588588588588589	0.00266710926726727\\
-0.585585585585586	0.0026685804954955\\
-0.582582582582583	0.00267005172372372\\
-0.57957957957958	0.00267152295195195\\
-0.576576576576577	0.00267299418018018\\
-0.573573573573574	0.00267446540840841\\
-0.570570570570571	0.00267593663663664\\
-0.567567567567568	0.00267740786486486\\
-0.564564564564564	0.00267887909309309\\
-0.561561561561562	0.00268035032132132\\
-0.558558558558559	0.00268182154954955\\
-0.555555555555556	0.00268329277777778\\
-0.552552552552553	0.00268476400600601\\
-0.549549549549549	0.00268623523423423\\
-0.546546546546547	0.00268770646246246\\
-0.543543543543544	0.00268917769069069\\
-0.540540540540541	0.00269064891891892\\
-0.537537537537538	0.00269212014714715\\
-0.534534534534534	0.00269359137537538\\
-0.531531531531532	0.0026950626036036\\
-0.528528528528529	0.00269653383183183\\
-0.525525525525526	0.00269800506006006\\
-0.522522522522523	0.00269947628828829\\
-0.519519519519519	0.00270094751651652\\
-0.516516516516517	0.00270241874474474\\
-0.513513513513513	0.00270388997297297\\
-0.510510510510511	0.0027053612012012\\
-0.507507507507508	0.00270683242942943\\
-0.504504504504504	0.00270830365765766\\
-0.501501501501502	0.00270977488588589\\
-0.498498498498498	0.00271124611411411\\
-0.495495495495496	0.00271271734234234\\
-0.492492492492492	0.00271418857057057\\
-0.489489489489489	0.0027156597987988\\
-0.486486486486487	0.00271713102702703\\
-0.483483483483483	0.00271860225525526\\
-0.480480480480481	0.00272007348348348\\
-0.477477477477477	0.00272154471171171\\
-0.474474474474474	0.00272301593993994\\
-0.471471471471472	0.00272448716816817\\
-0.468468468468468	0.0027259583963964\\
-0.465465465465465	0.00272742962462462\\
-0.462462462462462	0.00272890085285285\\
-0.459459459459459	0.00273037208108108\\
-0.456456456456456	0.00273184330930931\\
-0.453453453453453	0.00273331453753754\\
-0.45045045045045	0.00273478576576577\\
-0.447447447447447	0.00273625699399399\\
-0.444444444444444	0.00273772822222222\\
-0.441441441441441	0.00273919945045045\\
-0.438438438438438	0.00274067067867868\\
-0.435435435435435	0.00274214190690691\\
-0.432432432432432	0.00274361313513514\\
-0.429429429429429	0.00274508436336336\\
-0.426426426426426	0.00274655559159159\\
-0.423423423423423	0.00274802681981982\\
-0.42042042042042	0.00274949804804805\\
-0.417417417417417	0.00275096927627628\\
-0.414414414414414	0.0027524405045045\\
-0.411411411411411	0.00275391173273273\\
-0.408408408408408	0.00275538296096096\\
-0.405405405405405	0.00275685418918919\\
-0.402402402402402	0.00275832541741742\\
-0.399399399399399	0.00275979664564565\\
-0.396396396396396	0.00276126787387387\\
-0.393393393393393	0.0027627391021021\\
-0.39039039039039	0.00276421033033033\\
-0.387387387387387	0.00276568155855856\\
-0.384384384384384	0.00276715278678679\\
-0.381381381381381	0.00276862401501502\\
-0.378378378378378	0.00277009524324324\\
-0.375375375375375	0.00277156647147147\\
-0.372372372372372	0.0027730376996997\\
-0.369369369369369	0.00277450892792793\\
-0.366366366366366	0.00277598015615616\\
-0.363363363363363	0.00277745138438438\\
-0.36036036036036	0.00277892261261261\\
-0.357357357357357	0.00278039384084084\\
-0.354354354354354	0.00278186506906907\\
-0.351351351351351	0.0027833362972973\\
-0.348348348348348	0.00278480752552553\\
-0.345345345345345	0.00278627875375375\\
-0.342342342342342	0.00278774998198198\\
-0.339339339339339	0.00278922121021021\\
-0.336336336336336	0.00279069243843844\\
-0.333333333333333	0.00279216366666667\\
-0.33033033033033	0.00279363489489489\\
-0.327327327327327	0.00279510612312312\\
-0.324324324324324	0.00279657735135135\\
-0.321321321321321	0.00279804857957958\\
-0.318318318318318	0.00279951980780781\\
-0.315315315315315	0.00280099103603604\\
-0.312312312312312	0.00280246226426426\\
-0.309309309309309	0.00280393349249249\\
-0.306306306306306	0.00280540472072072\\
-0.303303303303303	0.00280687594894895\\
-0.3003003003003	0.00280834717717718\\
-0.297297297297297	0.00280981840540541\\
-0.294294294294294	0.00281128963363363\\
-0.291291291291291	0.00281276086186186\\
-0.288288288288288	0.00281423209009009\\
-0.285285285285285	0.00281570331831832\\
-0.282282282282282	0.00281717454654655\\
-0.279279279279279	0.00281864577477477\\
-0.276276276276276	0.002820117003003\\
-0.273273273273273	0.00282158823123123\\
-0.27027027027027	0.00282305945945946\\
-0.267267267267267	0.00282453068768769\\
-0.264264264264264	0.00282600191591592\\
-0.261261261261261	0.00282747314414414\\
-0.258258258258258	0.00282894437237237\\
-0.255255255255255	0.0028304156006006\\
-0.252252252252252	0.00283188682882883\\
-0.249249249249249	0.00283335805705706\\
-0.246246246246246	0.00283482928528529\\
-0.243243243243243	0.00283630051351351\\
-0.24024024024024	0.00283777174174174\\
-0.237237237237237	0.00283924296996997\\
-0.234234234234234	0.0028407141981982\\
-0.231231231231231	0.00284218542642643\\
-0.228228228228228	0.00284365665465465\\
-0.225225225225225	0.00284512788288288\\
-0.222222222222222	0.00284659911111111\\
-0.219219219219219	0.00284807033933934\\
-0.216216216216216	0.00284954156756757\\
-0.213213213213213	0.0028510127957958\\
-0.21021021021021	0.00285248402402402\\
-0.207207207207207	0.00285395525225225\\
-0.204204204204204	0.00285542648048048\\
-0.201201201201201	0.00285689770870871\\
-0.198198198198198	0.00285836893693694\\
-0.195195195195195	0.00285984016516517\\
-0.192192192192192	0.00286131139339339\\
-0.189189189189189	0.00286278262162162\\
-0.186186186186186	0.00286425384984985\\
-0.183183183183183	0.00286572507807808\\
-0.18018018018018	0.00286719630630631\\
-0.177177177177177	0.00286866753453453\\
-0.174174174174174	0.00287013876276276\\
-0.171171171171171	0.00287160999099099\\
-0.168168168168168	0.00287308121921922\\
-0.165165165165165	0.00287455244744745\\
-0.162162162162162	0.00287602367567568\\
-0.159159159159159	0.0028774949039039\\
-0.156156156156156	0.00287896613213213\\
-0.153153153153153	0.00288043736036036\\
-0.15015015015015	0.00288190858858859\\
-0.147147147147147	0.00288337981681682\\
-0.144144144144144	0.00288485104504504\\
-0.141141141141141	0.00288632227327327\\
-0.138138138138138	0.0028877935015015\\
-0.135135135135135	0.00288926472972973\\
-0.132132132132132	0.00289073595795796\\
-0.129129129129129	0.00289220718618619\\
-0.126126126126126	0.00289367841441441\\
-0.123123123123123	0.00289514964264264\\
-0.12012012012012	0.00289662087087087\\
-0.117117117117117	0.0028980920990991\\
-0.114114114114114	0.00289956332732733\\
-0.111111111111111	0.00290103455555556\\
-0.108108108108108	0.00290250578378378\\
-0.105105105105105	0.00290397701201201\\
-0.102102102102102	0.00290544824024024\\
-0.0990990990990991	0.00290691946846847\\
-0.0960960960960962	0.0029083906966967\\
-0.0930930930930931	0.00290986192492492\\
-0.0900900900900901	0.00291133315315315\\
-0.0870870870870871	0.00291280438138138\\
-0.0840840840840841	0.00291427560960961\\
-0.081081081081081	0.00291574683783784\\
-0.0780780780780781	0.00291721806606607\\
-0.075075075075075	0.00291868929429429\\
-0.0720720720720721	0.00292016052252252\\
-0.0690690690690691	0.00292163175075075\\
-0.066066066066066	0.00292310297897898\\
-0.0630630630630631	0.00292457420720721\\
-0.06006006006006	0.00292604543543544\\
-0.0570570570570571	0.00292751666366366\\
-0.0540540540540541	0.00292898789189189\\
-0.051051051051051	0.00293045912012012\\
-0.0480480480480481	0.00293193034834835\\
-0.045045045045045	0.00293340157657658\\
-0.0420420420420421	0.0029348728048048\\
-0.039039039039039	0.00293634403303303\\
-0.036036036036036	0.00293781526126126\\
-0.0330330330330331	0.00293928648948949\\
-0.03003003003003	0.00294075771771772\\
-0.027027027027027	0.00294222894594595\\
-0.024024024024024	0.00294370017417417\\
-0.021021021021021	0.0029451714024024\\
-0.0180180180180181	0.00294664263063063\\
-0.015015015015015	0.00294811385885886\\
-0.012012012012012	0.00294958508708709\\
-0.00900900900900903	0.00295105631531532\\
-0.00600600600600598	0.00295252754354354\\
-0.00300300300300305	0.00295399877177177\\
0	0.00295547\\
0.00300300300300305	0.00295694122822823\\
0.00600600600600609	0.00295841245645646\\
0.00900900900900892	0.00295988368468468\\
0.012012012012012	0.00296135491291291\\
0.015015015015015	0.00296282614114114\\
0.0180180180180181	0.00296429736936937\\
0.0210210210210211	0.0029657685975976\\
0.0240240240240239	0.00296723982582583\\
0.027027027027027	0.00296871105405405\\
0.03003003003003	0.00297018228228228\\
0.0330330330330331	0.00297165351051051\\
0.0360360360360361	0.00297312473873874\\
0.0390390390390389	0.00297459596696697\\
0.042042042042042	0.0029760671951952\\
0.045045045045045	0.00297753842342342\\
0.0480480480480481	0.00297900965165165\\
0.0510510510510511	0.00298048087987988\\
0.0540540540540539	0.00298195210810811\\
0.057057057057057	0.00298342333633634\\
0.06006006006006	0.00298489456456456\\
0.0630630630630631	0.00298636579279279\\
0.0660660660660661	0.00298783702102102\\
0.0690690690690692	0.00298930824924925\\
0.072072072072072	0.00299077947747748\\
0.075075075075075	0.00299225070570571\\
0.0780780780780781	0.00299372193393393\\
0.0810810810810811	0.00299519316216216\\
0.0840840840840842	0.00299666439039039\\
0.087087087087087	0.00299813561861862\\
0.0900900900900901	0.00299960684684685\\
0.0930930930930931	0.00300107807507507\\
0.0960960960960962	0.0030025493033033\\
0.0990990990990992	0.00300402053153153\\
0.102102102102102	0.00300549175975976\\
0.105105105105105	0.00300696298798799\\
0.108108108108108	0.00300843421621622\\
0.111111111111111	0.00300990544444444\\
0.114114114114114	0.00301137667267267\\
0.117117117117117	0.0030128479009009\\
0.12012012012012	0.00301431912912913\\
0.123123123123123	0.00301579035735736\\
0.126126126126126	0.00301726158558559\\
0.129129129129129	0.00301873281381381\\
0.132132132132132	0.00302020404204204\\
0.135135135135135	0.00302167527027027\\
0.138138138138138	0.0030231464984985\\
0.141141141141141	0.00302461772672673\\
0.144144144144144	0.00302608895495496\\
0.147147147147147	0.00302756018318318\\
0.15015015015015	0.00302903141141141\\
0.153153153153153	0.00303050263963964\\
0.156156156156156	0.00303197386786787\\
0.159159159159159	0.0030334450960961\\
0.162162162162162	0.00303491632432432\\
0.165165165165165	0.00303638755255255\\
0.168168168168168	0.00303785878078078\\
0.171171171171171	0.00303933000900901\\
0.174174174174174	0.00304080123723724\\
0.177177177177177	0.00304227246546547\\
0.18018018018018	0.00304374369369369\\
0.183183183183183	0.00304521492192192\\
0.186186186186186	0.00304668615015015\\
0.189189189189189	0.00304815737837838\\
0.192192192192192	0.00304962860660661\\
0.195195195195195	0.00305109983483483\\
0.198198198198198	0.00305257106306306\\
0.201201201201201	0.00305404229129129\\
0.204204204204204	0.00305551351951952\\
0.207207207207207	0.00305698474774775\\
0.21021021021021	0.00305845597597598\\
0.213213213213213	0.0030599272042042\\
0.216216216216216	0.00306139843243243\\
0.219219219219219	0.00306286966066066\\
0.222222222222222	0.00306434088888889\\
0.225225225225225	0.00306581211711712\\
0.228228228228228	0.00306728334534535\\
0.231231231231231	0.00306875457357357\\
0.234234234234234	0.0030702258018018\\
0.237237237237237	0.00307169703003003\\
0.24024024024024	0.00307316825825826\\
0.243243243243243	0.00307463948648649\\
0.246246246246246	0.00307611071471471\\
0.249249249249249	0.00307758194294294\\
0.252252252252252	0.00307905317117117\\
0.255255255255255	0.0030805243993994\\
0.258258258258258	0.00308199562762763\\
0.261261261261261	0.00308346685585586\\
0.264264264264264	0.00308493808408408\\
0.267267267267267	0.00308640931231231\\
0.27027027027027	0.00308788054054054\\
0.273273273273273	0.00308935176876877\\
0.276276276276276	0.003090822996997\\
0.279279279279279	0.00309229422522523\\
0.282282282282282	0.00309376545345345\\
0.285285285285285	0.00309523668168168\\
0.288288288288288	0.00309670790990991\\
0.291291291291291	0.00309817913813814\\
0.294294294294294	0.00309965036636637\\
0.297297297297297	0.00310112159459459\\
0.3003003003003	0.00310259282282282\\
0.303303303303303	0.00310406405105105\\
0.306306306306306	0.00310553527927928\\
0.309309309309309	0.00310700650750751\\
0.312312312312312	0.00310847773573574\\
0.315315315315315	0.00310994896396396\\
0.318318318318318	0.00311142019219219\\
0.321321321321321	0.00311289142042042\\
0.324324324324324	0.00311436264864865\\
0.327327327327327	0.00311583387687688\\
0.33033033033033	0.00311730510510511\\
0.333333333333333	0.00311877633333333\\
0.336336336336336	0.00312024756156156\\
0.339339339339339	0.00312171878978979\\
0.342342342342342	0.00312319001801802\\
0.345345345345345	0.00312466124624625\\
0.348348348348348	0.00312613247447447\\
0.351351351351351	0.0031276037027027\\
0.354354354354354	0.00312907493093093\\
0.357357357357357	0.00313054615915916\\
0.36036036036036	0.00313201738738739\\
0.363363363363363	0.00313348861561562\\
0.366366366366366	0.00313495984384384\\
0.369369369369369	0.00313643107207207\\
0.372372372372372	0.0031379023003003\\
0.375375375375375	0.00313937352852853\\
0.378378378378378	0.00314084475675676\\
0.381381381381381	0.00314231598498498\\
0.384384384384384	0.00314378721321321\\
0.387387387387387	0.00314525844144144\\
0.39039039039039	0.00314672966966967\\
0.393393393393393	0.0031482008978979\\
0.396396396396396	0.00314967212612613\\
0.399399399399399	0.00315114335435435\\
0.402402402402402	0.00315261458258258\\
0.405405405405405	0.00315408581081081\\
0.408408408408408	0.00315555703903904\\
0.411411411411411	0.00315702826726727\\
0.414414414414414	0.0031584994954955\\
0.417417417417417	0.00315997072372372\\
0.42042042042042	0.00316144195195195\\
0.423423423423423	0.00316291318018018\\
0.426426426426426	0.00316438440840841\\
0.429429429429429	0.00316585563663664\\
0.432432432432432	0.00316732686486486\\
0.435435435435436	0.00316879809309309\\
0.438438438438439	0.00317026932132132\\
0.441441441441441	0.00317174054954955\\
0.444444444444444	0.00317321177777778\\
0.447447447447447	0.00317468300600601\\
0.450450450450451	0.00317615423423423\\
0.453453453453454	0.00317762546246246\\
0.456456456456456	0.00317909669069069\\
0.459459459459459	0.00318056791891892\\
0.462462462462462	0.00318203914714715\\
0.465465465465466	0.00318351037537538\\
0.468468468468469	0.0031849816036036\\
0.471471471471471	0.00318645283183183\\
0.474474474474474	0.00318792406006006\\
0.477477477477477	0.00318939528828829\\
0.480480480480481	0.00319086651651652\\
0.483483483483484	0.00319233774474474\\
0.486486486486486	0.00319380897297297\\
0.489489489489489	0.0031952802012012\\
0.492492492492492	0.00319675142942943\\
0.495495495495496	0.00319822265765766\\
0.498498498498499	0.00319969388588589\\
0.501501501501501	0.00320116511411411\\
0.504504504504504	0.00320263634234234\\
0.507507507507508	0.00320410757057057\\
0.510510510510511	0.0032055787987988\\
0.513513513513514	0.00320705002702703\\
0.516516516516516	0.00320852125525526\\
0.519519519519519	0.00320999248348348\\
0.522522522522523	0.00321146371171171\\
0.525525525525526	0.00321293493993994\\
0.528528528528529	0.00321440616816817\\
0.531531531531531	0.0032158773963964\\
0.534534534534534	0.00321734862462462\\
0.537537537537538	0.00321881985285285\\
0.540540540540541	0.00322029108108108\\
0.543543543543544	0.00322176230930931\\
0.546546546546546	0.00322323353753754\\
0.549549549549549	0.00322470476576577\\
0.552552552552553	0.00322617599399399\\
0.555555555555556	0.00322764722222222\\
0.558558558558559	0.00322911845045045\\
0.561561561561561	0.00323058967867868\\
0.564564564564564	0.00323206090690691\\
0.567567567567568	0.00323353213513514\\
0.570570570570571	0.00323500336336336\\
0.573573573573574	0.00323647459159159\\
0.576576576576577	0.00323794581981982\\
0.57957957957958	0.00323941704804805\\
0.582582582582583	0.00324088827627628\\
0.585585585585586	0.0032423595045045\\
0.588588588588589	0.00324383073273273\\
0.591591591591592	0.00324530196096096\\
0.594594594594595	0.00324677318918919\\
0.597597597597598	0.00324824441741742\\
0.600600600600601	0.00324971564564565\\
0.603603603603604	0.00325118687387387\\
0.606606606606607	0.0032526581021021\\
0.60960960960961	0.00325412933033033\\
0.612612612612613	0.00325560055855856\\
0.615615615615616	0.00325707178678679\\
0.618618618618619	0.00325854301501501\\
0.621621621621622	0.00326001424324324\\
0.624624624624625	0.00326148547147147\\
0.627627627627628	0.0032629566996997\\
0.630630630630631	0.00326442792792793\\
0.633633633633634	0.00326589915615616\\
0.636636636636637	0.00326737038438438\\
0.63963963963964	0.00326884161261261\\
0.642642642642643	0.00327031284084084\\
0.645645645645646	0.00327178406906907\\
0.648648648648649	0.0032732552972973\\
0.651651651651652	0.00327472652552553\\
0.654654654654655	0.00327619775375375\\
0.657657657657658	0.00327766898198198\\
0.660660660660661	0.00327914021021021\\
0.663663663663664	0.00328061143843844\\
0.666666666666667	0.00328208266666667\\
0.66966966966967	0.00328355389489489\\
0.672672672672673	0.00328502512312312\\
0.675675675675676	0.00328649635135135\\
0.678678678678679	0.00328796757957958\\
0.681681681681682	0.00328943880780781\\
0.684684684684685	0.00329091003603604\\
0.687687687687688	0.00329238126426426\\
0.690690690690691	0.00329385249249249\\
0.693693693693694	0.00329532372072072\\
0.696696696696697	0.00329679494894895\\
0.6996996996997	0.00329826617717718\\
0.702702702702703	0.00329973740540541\\
0.705705705705706	0.00330120863363363\\
0.708708708708709	0.00330267986186186\\
0.711711711711712	0.00330415109009009\\
0.714714714714715	0.00330562231831832\\
0.717717717717718	0.00330709354654655\\
0.720720720720721	0.00330856477477477\\
0.723723723723724	0.003310036003003\\
0.726726726726727	0.00331150723123123\\
0.72972972972973	0.00331297845945946\\
0.732732732732733	0.00331444968768769\\
0.735735735735736	0.00331592091591592\\
0.738738738738739	0.00331739214414414\\
0.741741741741742	0.00331886337237237\\
0.744744744744745	0.0033203346006006\\
0.747747747747748	0.00332180582882883\\
0.750750750750751	0.00332327705705706\\
0.753753753753754	0.00332474828528529\\
0.756756756756757	0.00332621951351351\\
0.75975975975976	0.00332769074174174\\
0.762762762762763	0.00332916196996997\\
0.765765765765766	0.0033306331981982\\
0.768768768768769	0.00333210442642643\\
0.771771771771772	0.00333357565465465\\
0.774774774774775	0.00333504688288288\\
0.777777777777778	0.00333651811111111\\
0.780780780780781	0.00333798933933934\\
0.783783783783784	0.00333946056756757\\
0.786786786786787	0.0033409317957958\\
0.78978978978979	0.00334240302402402\\
0.792792792792793	0.00334387425225225\\
0.795795795795796	0.00334534548048048\\
0.798798798798799	0.00334681670870871\\
0.801801801801802	0.00334828793693694\\
0.804804804804805	0.00334975916516517\\
0.807807807807808	0.00335123039339339\\
0.810810810810811	0.00335270162162162\\
0.813813813813814	0.00335417284984985\\
0.816816816816817	0.00335564407807808\\
0.81981981981982	0.00335711530630631\\
0.822822822822823	0.00335858653453453\\
0.825825825825826	0.00336005776276276\\
0.828828828828829	0.00336152899099099\\
0.831831831831832	0.00336300021921922\\
0.834834834834835	0.00336447144744745\\
0.837837837837838	0.00336594267567568\\
0.840840840840841	0.0033674139039039\\
0.843843843843844	0.00336888513213213\\
0.846846846846847	0.00337035636036036\\
0.84984984984985	0.00337182758858859\\
0.852852852852853	0.00337329881681682\\
0.855855855855856	0.00337477004504504\\
0.858858858858859	0.00337624127327327\\
0.861861861861862	0.0033777125015015\\
0.864864864864865	0.00337918372972973\\
0.867867867867868	0.00338065495795796\\
0.870870870870871	0.00338212618618619\\
0.873873873873874	0.00338359741441441\\
0.876876876876877	0.00338506864264264\\
0.87987987987988	0.00338653987087087\\
0.882882882882883	0.0033880110990991\\
0.885885885885886	0.00338948232732733\\
0.888888888888889	0.00339095355555556\\
0.891891891891892	0.00339242478378378\\
0.894894894894895	0.00339389601201201\\
0.897897897897898	0.00339536724024024\\
0.900900900900901	0.00339683846846847\\
0.903903903903904	0.0033983096966967\\
0.906906906906907	0.00339978092492492\\
0.90990990990991	0.00340125215315315\\
0.912912912912913	0.00340272338138138\\
0.915915915915916	0.00340419460960961\\
0.918918918918919	0.00340566583783784\\
0.921921921921922	0.00340713706606607\\
0.924924924924925	0.00340860829429429\\
0.927927927927928	0.00341007952252252\\
0.930930930930931	0.00341155075075075\\
0.933933933933934	0.00341302197897898\\
0.936936936936937	0.00341449320720721\\
0.93993993993994	0.00341596443543544\\
0.942942942942943	0.00341743566366366\\
0.945945945945946	0.00341890689189189\\
0.948948948948949	0.00342037812012012\\
0.951951951951952	0.00342184934834835\\
0.954954954954955	0.00342332057657658\\
0.957957957957958	0.0034247918048048\\
0.960960960960961	0.00342626303303303\\
0.963963963963964	0.00342773426126126\\
0.966966966966967	0.00342920548948949\\
0.96996996996997	0.00343067671771772\\
0.972972972972973	0.00343214794594595\\
0.975975975975976	0.00343361917417417\\
0.978978978978979	0.0034350904024024\\
0.981981981981982	0.00343656163063063\\
0.984984984984985	0.00343803285885886\\
0.987987987987988	0.00343950408708709\\
0.990990990990991	0.00344097531531532\\
0.993993993993994	0.00344244654354354\\
0.996996996996997	0.00344391777177177\\
1	0.003445389\\
1.003003003003	0.00344686022822823\\
1.00600600600601	0.00344833145645646\\
1.00900900900901	0.00344980268468468\\
1.01201201201201	0.00345127391291291\\
1.01501501501502	0.00345274514114114\\
1.01801801801802	0.00345421636936937\\
1.02102102102102	0.0034556875975976\\
1.02402402402402	0.00345715882582583\\
1.02702702702703	0.00345863005405405\\
1.03003003003003	0.00346010128228228\\
1.03303303303303	0.00346157251051051\\
1.03603603603604	0.00346304373873874\\
1.03903903903904	0.00346451496696697\\
1.04204204204204	0.0034659861951952\\
1.04504504504505	0.00346745742342342\\
1.04804804804805	0.00346892865165165\\
1.05105105105105	0.00347039987987988\\
1.05405405405405	0.00347187110810811\\
1.05705705705706	0.00347334233633634\\
1.06006006006006	0.00347481356456456\\
1.06306306306306	0.00347628479279279\\
1.06606606606607	0.00347775602102102\\
1.06906906906907	0.00347922724924925\\
1.07207207207207	0.00348069847747748\\
1.07507507507508	0.00348216970570571\\
1.07807807807808	0.00348364093393393\\
1.08108108108108	0.00348511216216216\\
1.08408408408408	0.00348658339039039\\
1.08708708708709	0.00348805461861862\\
1.09009009009009	0.00348952584684685\\
1.09309309309309	0.00349099707507508\\
1.0960960960961	0.0034924683033033\\
1.0990990990991	0.00349393953153153\\
1.1021021021021	0.00349541075975976\\
1.10510510510511	0.00349688198798799\\
1.10810810810811	0.00349835321621622\\
1.11111111111111	0.00349982444444444\\
1.11411411411411	0.00350129567267267\\
1.11711711711712	0.0035027669009009\\
1.12012012012012	0.00350423812912913\\
1.12312312312312	0.00350570935735736\\
1.12612612612613	0.00350718058558559\\
1.12912912912913	0.00350865181381381\\
1.13213213213213	0.00351012304204204\\
1.13513513513514	0.00351159427027027\\
1.13813813813814	0.0035130654984985\\
1.14114114114114	0.00351453672672673\\
1.14414414414414	0.00351600795495495\\
1.14714714714715	0.00351747918318318\\
1.15015015015015	0.00351895041141141\\
1.15315315315315	0.00352042163963964\\
1.15615615615616	0.00352189286786787\\
1.15915915915916	0.0035233640960961\\
1.16216216216216	0.00352483532432432\\
1.16516516516517	0.00352630655255255\\
1.16816816816817	0.00352777778078078\\
1.17117117117117	0.00352924900900901\\
1.17417417417417	0.00353072023723724\\
1.17717717717718	0.00353219146546547\\
1.18018018018018	0.00353366269369369\\
1.18318318318318	0.00353513392192192\\
1.18618618618619	0.00353660515015015\\
1.18918918918919	0.00353807637837838\\
1.19219219219219	0.00353954760660661\\
1.1951951951952	0.00354101883483483\\
1.1981981981982	0.00354249006306306\\
1.2012012012012	0.00354396129129129\\
1.2042042042042	0.00354543251951952\\
1.20720720720721	0.00354690374774775\\
1.21021021021021	0.00354837497597598\\
1.21321321321321	0.0035498462042042\\
1.21621621621622	0.00355131743243243\\
1.21921921921922	0.00355278866066066\\
1.22222222222222	0.00355425988888889\\
1.22522522522523	0.00355573111711712\\
1.22822822822823	0.00355720234534535\\
1.23123123123123	0.00355867357357357\\
1.23423423423423	0.0035601448018018\\
1.23723723723724	0.00356161603003003\\
1.24024024024024	0.00356308725825826\\
1.24324324324324	0.00356455848648649\\
1.24624624624625	0.00356602971471471\\
1.24924924924925	0.00356750094294294\\
1.25225225225225	0.00356897217117117\\
1.25525525525526	0.0035704433993994\\
1.25825825825826	0.00357191462762763\\
1.26126126126126	0.00357338585585586\\
1.26426426426426	0.00357485708408408\\
1.26726726726727	0.00357632831231231\\
1.27027027027027	0.00357779954054054\\
1.27327327327327	0.00357927076876877\\
1.27627627627628	0.003580741996997\\
1.27927927927928	0.00358221322522523\\
1.28228228228228	0.00358368445345345\\
1.28528528528529	0.00358515568168168\\
1.28828828828829	0.00358662690990991\\
1.29129129129129	0.00358809813813814\\
1.29429429429429	0.00358956936636637\\
1.2972972972973	0.00359104059459459\\
1.3003003003003	0.00359251182282282\\
1.3033033033033	0.00359398305105105\\
1.30630630630631	0.00359545427927928\\
1.30930930930931	0.00359692550750751\\
1.31231231231231	0.00359839673573574\\
1.31531531531532	0.00359986796396396\\
1.31831831831832	0.00360133919219219\\
1.32132132132132	0.00360281042042042\\
1.32432432432432	0.00360428164864865\\
1.32732732732733	0.00360575287687688\\
1.33033033033033	0.00360722410510511\\
1.33333333333333	0.00360869533333333\\
1.33633633633634	0.00361016656156156\\
1.33933933933934	0.00361163778978979\\
1.34234234234234	0.00361310901801802\\
1.34534534534535	0.00361458024624625\\
1.34834834834835	0.00361605147447447\\
1.35135135135135	0.0036175227027027\\
1.35435435435435	0.00361899393093093\\
1.35735735735736	0.00362046515915916\\
1.36036036036036	0.00362193638738739\\
1.36336336336336	0.00362340761561562\\
1.36636636636637	0.00362487884384384\\
1.36936936936937	0.00362635007207207\\
1.37237237237237	0.0036278213003003\\
1.37537537537538	0.00362929252852853\\
1.37837837837838	0.00363076375675676\\
1.38138138138138	0.00363223498498498\\
1.38438438438438	0.00363370621321321\\
1.38738738738739	0.00363517744144144\\
1.39039039039039	0.00363664866966967\\
1.39339339339339	0.0036381198978979\\
1.3963963963964	0.00363959112612613\\
1.3993993993994	0.00364106235435435\\
1.4024024024024	0.00364253358258258\\
1.40540540540541	0.00364400481081081\\
1.40840840840841	0.00364547603903904\\
1.41141141141141	0.00364694726726727\\
1.41441441441441	0.0036484184954955\\
1.41741741741742	0.00364988972372372\\
1.42042042042042	0.00365136095195195\\
1.42342342342342	0.00365283218018018\\
1.42642642642643	0.00365430340840841\\
1.42942942942943	0.00365577463663664\\
1.43243243243243	0.00365724586486486\\
1.43543543543544	0.00365871709309309\\
1.43843843843844	0.00366018832132132\\
1.44144144144144	0.00366165954954955\\
1.44444444444444	0.00366313077777778\\
1.44744744744745	0.00366460200600601\\
1.45045045045045	0.00366607323423423\\
1.45345345345345	0.00366754446246246\\
1.45645645645646	0.00366901569069069\\
1.45945945945946	0.00367048691891892\\
1.46246246246246	0.00367195814714715\\
1.46546546546547	0.00367342937537538\\
1.46846846846847	0.0036749006036036\\
1.47147147147147	0.00367637183183183\\
1.47447447447447	0.00367784306006006\\
1.47747747747748	0.00367931428828829\\
1.48048048048048	0.00368078551651652\\
1.48348348348348	0.00368225674474474\\
1.48648648648649	0.00368372797297297\\
1.48948948948949	0.0036851992012012\\
1.49249249249249	0.00368667042942943\\
1.4954954954955	0.00368814165765766\\
1.4984984984985	0.00368961288588589\\
1.5015015015015	0.00369108411411411\\
1.5045045045045	0.00369255534234234\\
1.50750750750751	0.00369402657057057\\
1.51051051051051	0.0036954977987988\\
1.51351351351351	0.00369696902702703\\
1.51651651651652	0.00369844025525526\\
1.51951951951952	0.00369991148348348\\
1.52252252252252	0.00370138271171171\\
1.52552552552553	0.00370285393993994\\
1.52852852852853	0.00370432516816817\\
1.53153153153153	0.0037057963963964\\
1.53453453453453	0.00370726762462462\\
1.53753753753754	0.00370873885285285\\
1.54054054054054	0.00371021008108108\\
1.54354354354354	0.00371168130930931\\
1.54654654654655	0.00371315253753754\\
1.54954954954955	0.00371462376576577\\
1.55255255255255	0.00371609499399399\\
1.55555555555556	0.00371756622222222\\
1.55855855855856	0.00371903745045045\\
1.56156156156156	0.00372050867867868\\
1.56456456456456	0.00372197990690691\\
1.56756756756757	0.00372345113513513\\
1.57057057057057	0.00372492236336336\\
1.57357357357357	0.00372639359159159\\
1.57657657657658	0.00372786481981982\\
1.57957957957958	0.00372933604804805\\
1.58258258258258	0.00373080727627628\\
1.58558558558559	0.0037322785045045\\
1.58858858858859	0.00373374973273273\\
1.59159159159159	0.00373522096096096\\
1.59459459459459	0.00373669218918919\\
1.5975975975976	0.00373816341741742\\
1.6006006006006	0.00373963464564565\\
1.6036036036036	0.00374110587387387\\
1.60660660660661	0.0037425771021021\\
1.60960960960961	0.00374404833033033\\
1.61261261261261	0.00374551955855856\\
1.61561561561562	0.00374699078678679\\
1.61861861861862	0.00374846201501501\\
1.62162162162162	0.00374993324324324\\
1.62462462462462	0.00375140447147147\\
1.62762762762763	0.0037528756996997\\
1.63063063063063	0.00375434692792793\\
1.63363363363363	0.00375581815615616\\
1.63663663663664	0.00375728938438438\\
1.63963963963964	0.00375876061261261\\
1.64264264264264	0.00376023184084084\\
1.64564564564565	0.00376170306906907\\
1.64864864864865	0.0037631742972973\\
1.65165165165165	0.00376464552552553\\
1.65465465465465	0.00376611675375375\\
1.65765765765766	0.00376758798198198\\
1.66066066066066	0.00376905921021021\\
1.66366366366366	0.00377053043843844\\
1.66666666666667	0.00377200166666667\\
1.66966966966967	0.0037734728948949\\
1.67267267267267	0.00377494412312312\\
1.67567567567568	0.00377641535135135\\
1.67867867867868	0.00377788657957958\\
1.68168168168168	0.00377935780780781\\
1.68468468468468	0.00378082903603604\\
1.68768768768769	0.00378230026426426\\
1.69069069069069	0.00378377149249249\\
1.69369369369369	0.00378524272072072\\
1.6966966966967	0.00378671394894895\\
1.6996996996997	0.00378818517717718\\
1.7027027027027	0.00378965640540541\\
1.70570570570571	0.00379112763363363\\
1.70870870870871	0.00379259886186186\\
1.71171171171171	0.00379407009009009\\
1.71471471471471	0.00379554131831832\\
1.71771771771772	0.00379701254654655\\
1.72072072072072	0.00379848377477477\\
1.72372372372372	0.003799955003003\\
1.72672672672673	0.00380142623123123\\
1.72972972972973	0.00380289745945946\\
1.73273273273273	0.00380436868768769\\
1.73573573573574	0.00380583991591592\\
1.73873873873874	0.00380731114414414\\
1.74174174174174	0.00380878237237237\\
1.74474474474474	0.0038102536006006\\
1.74774774774775	0.00381172482882883\\
1.75075075075075	0.00381319605705706\\
1.75375375375375	0.00381466728528529\\
1.75675675675676	0.00381613851351351\\
1.75975975975976	0.00381760974174174\\
1.76276276276276	0.00381908096996997\\
1.76576576576577	0.0038205521981982\\
1.76876876876877	0.00382202342642643\\
1.77177177177177	0.00382349465465465\\
1.77477477477477	0.00382496588288288\\
1.77777777777778	0.00382643711111111\\
1.78078078078078	0.00382790833933934\\
1.78378378378378	0.00382937956756757\\
1.78678678678679	0.0038308507957958\\
1.78978978978979	0.00383232202402402\\
1.79279279279279	0.00383379325225225\\
1.7957957957958	0.00383526448048048\\
1.7987987987988	0.00383673570870871\\
1.8018018018018	0.00383820693693694\\
1.8048048048048	0.00383967816516517\\
1.80780780780781	0.00384114939339339\\
1.81081081081081	0.00384262062162162\\
1.81381381381381	0.00384409184984985\\
1.81681681681682	0.00384556307807808\\
1.81981981981982	0.00384703430630631\\
1.82282282282282	0.00384850553453453\\
1.82582582582583	0.00384997676276276\\
1.82882882882883	0.00385144799099099\\
1.83183183183183	0.00385291921921922\\
1.83483483483483	0.00385439044744745\\
1.83783783783784	0.00385586167567568\\
1.84084084084084	0.0038573329039039\\
1.84384384384384	0.00385880413213213\\
1.84684684684685	0.00386027536036036\\
1.84984984984985	0.00386174658858859\\
1.85285285285285	0.00386321781681682\\
1.85585585585586	0.00386468904504505\\
1.85885885885886	0.00386616027327327\\
1.86186186186186	0.0038676315015015\\
1.86486486486486	0.00386910272972973\\
1.86786786786787	0.00387057395795796\\
1.87087087087087	0.00387204518618619\\
1.87387387387387	0.00387351641441441\\
1.87687687687688	0.00387498764264264\\
1.87987987987988	0.00387645887087087\\
1.88288288288288	0.0038779300990991\\
1.88588588588589	0.00387940132732733\\
1.88888888888889	0.00388087255555556\\
1.89189189189189	0.00388234378378378\\
1.89489489489489	0.00388381501201201\\
1.8978978978979	0.00388528624024024\\
1.9009009009009	0.00388675746846847\\
1.9039039039039	0.0038882286966967\\
1.90690690690691	0.00388969992492492\\
1.90990990990991	0.00389117115315315\\
1.91291291291291	0.00389264238138138\\
1.91591591591592	0.00389411360960961\\
1.91891891891892	0.00389558483783784\\
1.92192192192192	0.00389705606606607\\
1.92492492492492	0.00389852729429429\\
1.92792792792793	0.00389999852252252\\
1.93093093093093	0.00390146975075075\\
1.93393393393393	0.00390294097897898\\
1.93693693693694	0.00390441220720721\\
1.93993993993994	0.00390588343543544\\
1.94294294294294	0.00390735466366366\\
1.94594594594595	0.00390882589189189\\
1.94894894894895	0.00391029712012012\\
1.95195195195195	0.00391176834834835\\
1.95495495495495	0.00391323957657658\\
1.95795795795796	0.0039147108048048\\
1.96096096096096	0.00391618203303303\\
1.96396396396396	0.00391765326126126\\
1.96696696696697	0.00391912448948949\\
1.96996996996997	0.00392059571771772\\
1.97297297297297	0.00392206694594595\\
1.97597597597598	0.00392353817417417\\
1.97897897897898	0.0039250094024024\\
1.98198198198198	0.00392648063063063\\
1.98498498498498	0.00392795185885886\\
1.98798798798799	0.00392942308708709\\
1.99099099099099	0.00393089431531532\\
1.99399399399399	0.00393236554354354\\
1.996996996997	0.00393383677177177\\
2	0.003935308\\
};
\addlegendentry{Linear}

\end{axis}
\end{tikzpicture}%
\end{figure}
For the heating curve we get the following regression parameters

\noindent For cubic fit:\begin{align*}
    A &= 2.95208 \times 10^{-3}\\
    B &=0.60321 \times 10^{-3}\\
    C &=-0.159792 \times 10^{-3}\\
    R^{2}&=0.9927
\end{align*}
For linear fit:
\begin{align*}
    a_0 &= 2.95547 \times 10^{-3} \\
    a_1 &=0.489919 \times 10^{-3} \\
    R^{2}&=0.9769
\end{align*}

\begin{figure}[H]
    \centering
    \setlength\fwidth{0.7\textwidth}
    \setlength\fheight{0.5\textwidth}
    % This file was created by matlab2tikz.
%
\definecolor{mycolor1}{rgb}{0.00000,0.44700,0.74100}%
\definecolor{mycolor2}{rgb}{0.85000,0.32500,0.09800}%
\definecolor{mycolor3}{rgb}{0.12941,0.12941,0.12941}%
%
\begin{tikzpicture}

\begin{axis}[%
width=0.951\fwidth,
height=\fheight,
at={(0\fwidth,0\fheight)},
scale only axis,
xmin=-0.5,
xmax=1.35,
xlabel style={font=\color{mycolor3}},
xlabel={ln(R)},
ymin=0.0024,
ymax=0.004,
ylabel style={font=\color{mycolor3}},
ylabel={1/T},
axis background/.style={fill=white},
title style={font=\bfseries\color{mycolor3}},
title={Cooling},
xmajorgrids,
ymajorgrids,
legend style={at={(0.03,0.97)}, anchor=north west, legend cell align=left, align=left}
]
\addplot[only marks, mark=o, mark options={}, mark size=1.5000pt, draw=mycolor1] table[row sep=crcr]{%
x	y\\
1.03462845662532	0.00341296928327645\\
1.03462845662532	0.00337837837837838\\
1.01494034979294	0.00335570469798658\\
0.974050902792877	0.00332225913621262\\
0.943494515906102	0.0033003300330033\\
0.870988813760575	0.00324675324675325\\
0.800717078282385	0.00319488817891374\\
0.726727209026572	0.00314465408805031\\
0.655138434811382	0.00309597523219814\\
0.551449997972875	0.00304878048780488\\
0.444044795918076	0.003003003003003\\
0.385606273598312	0.0029585798816568\\
0.342422680822206	0.00291545189504373\\
0.309630167425899	0.0028735632183908\\
0.184691430817599	0.0028328611898017\\
-0.168130225719498	0.00280898876404494\\
};
\addlegendentry{Data}

\addplot [color=mycolor2]
  table[row sep=crcr]{%
-1	0.002287756\\
-0.996996996996997	0.00229077358442982\\
-0.993993993993994	0.00229377843536061\\
-0.990990990990991	0.0022967705911463\\
-0.987987987987988	0.00229975009014079\\
-0.984984984984985	0.00230271697069799\\
-0.981981981981982	0.00230567127117183\\
-0.978978978978979	0.00230861302991621\\
-0.975975975975976	0.00231154228528504\\
-0.972972972972973	0.00231445907563224\\
-0.96996996996997	0.00231736343931172\\
-0.966966966966967	0.0023202554146774\\
-0.963963963963964	0.00232313504008318\\
-0.960960960960961	0.00232600235388298\\
-0.957957957957958	0.00232885739443071\\
-0.954954954954955	0.00233170020008028\\
-0.951951951951952	0.00233453080918562\\
-0.948948948948949	0.00233734926010062\\
-0.945945945945946	0.0023401555911792\\
-0.942942942942943	0.00234294984077528\\
-0.93993993993994	0.00234573204724276\\
-0.936936936936937	0.00234850224893557\\
-0.933933933933934	0.00235126048420761\\
-0.930930930930931	0.00235400679141279\\
-0.927927927927928	0.00235674120890503\\
-0.924924924924925	0.00235946377503825\\
-0.921921921921922	0.00236217452816635\\
-0.918918918918919	0.00236487350664324\\
-0.915915915915916	0.00236756074882284\\
-0.912912912912913	0.00237023629305907\\
-0.90990990990991	0.00237290017770583\\
-0.906906906906907	0.00237555244111704\\
-0.903903903903904	0.0023781931216466\\
-0.900900900900901	0.00238082225764844\\
-0.897897897897898	0.00238343988747647\\
-0.894894894894895	0.00238604604948459\\
-0.891891891891892	0.00238864078202673\\
-0.888888888888889	0.00239122412345679\\
-0.885885885885886	0.00239379611212869\\
-0.882882882882883	0.00239635678639633\\
-0.87987987987988	0.00239890618461364\\
-0.876876876876877	0.00240144434513452\\
-0.873873873873874	0.00240397130631289\\
-0.870870870870871	0.00240648710650266\\
-0.867867867867868	0.00240899178405774\\
-0.864864864864865	0.00241148537733204\\
-0.861861861861862	0.00241396792467949\\
-0.858858858858859	0.00241643946445398\\
-0.855855855855856	0.00241890003500944\\
-0.852852852852853	0.00242134967469978\\
-0.84984984984985	0.0024237884218789\\
-0.846846846846847	0.00242621631490073\\
-0.843843843843844	0.00242863339211917\\
-0.840840840840841	0.00243103969188814\\
-0.837837837837838	0.00243343525256155\\
-0.834834834834835	0.00243582011249331\\
-0.831831831831832	0.00243819431003733\\
-0.828828828828829	0.00244055788354754\\
-0.825825825825826	0.00244291087137783\\
-0.822822822822823	0.00244525331188213\\
-0.81981981981982	0.00244758524341434\\
-0.816816816816817	0.00244990670432838\\
-0.813813813813814	0.00245221773297817\\
-0.810810810810811	0.00245451836771761\\
-0.807807807807808	0.00245680864690061\\
-0.804804804804805	0.0024590886088811\\
-0.801801801801802	0.00246135829201298\\
-0.798798798798799	0.00246361773465016\\
-0.795795795795796	0.00246586697514656\\
-0.792792792792793	0.00246810605185609\\
-0.78978978978979	0.00247033500313267\\
-0.786786786786787	0.0024725538673302\\
-0.783783783783784	0.0024747626828026\\
-0.780780780780781	0.00247696148790378\\
-0.777777777777778	0.00247915032098765\\
-0.774774774774775	0.00248132922040814\\
-0.771771771771772	0.00248349822451914\\
-0.768768768768769	0.00248565737167457\\
-0.765765765765766	0.00248780670022835\\
-0.762762762762763	0.00248994624853439\\
-0.75975975975976	0.0024920760549466\\
-0.756756756756757	0.00249419615781889\\
-0.753753753753754	0.00249630659550517\\
-0.750750750750751	0.00249840740635937\\
-0.747747747747748	0.00250049862873538\\
-0.744744744744745	0.00250258030098713\\
-0.741741741741742	0.00250465246146853\\
-0.738738738738739	0.00250671514853349\\
-0.735735735735736	0.00250876840053591\\
-0.732732732732733	0.00251081225582973\\
-0.72972972972973	0.00251284675276884\\
-0.726726726726727	0.00251487192970716\\
-0.723723723723724	0.0025168878249986\\
-0.720720720720721	0.00251889447699708\\
-0.717717717717718	0.00252089192405651\\
-0.714714714714715	0.0025228802045308\\
-0.711711711711712	0.00252485935677387\\
-0.708708708708709	0.00252682941913962\\
-0.705705705705706	0.00252879042998197\\
-0.702702702702703	0.00253074242765483\\
-0.6996996996997	0.00253268545051211\\
-0.696696696696697	0.00253461953690774\\
-0.693693693693694	0.00253654472519561\\
-0.690690690690691	0.00253846105372965\\
-0.687687687687688	0.00254036856086376\\
-0.684684684684685	0.00254226728495186\\
-0.681681681681682	0.00254415726434786\\
-0.678678678678679	0.00254603853740568\\
-0.675675675675676	0.00254791114247922\\
-0.672672672672673	0.0025497751179224\\
-0.66966966966967	0.00255163050208914\\
-0.666666666666667	0.00255347733333333\\
-0.663663663663664	0.00255531565000891\\
-0.660660660660661	0.00255714549046977\\
-0.657657657657658	0.00255896689306984\\
-0.654654654654655	0.00256077989616302\\
-0.651651651651652	0.00256258453810324\\
-0.648648648648649	0.00256438085724439\\
-0.645645645645646	0.00256616889194039\\
-0.642642642642643	0.00256794868054517\\
-0.63963963963964	0.00256972026141262\\
-0.636636636636637	0.00257148367289666\\
-0.633633633633634	0.00257323895335121\\
-0.630630630630631	0.00257498614113017\\
-0.627627627627628	0.00257672527458747\\
-0.624624624624625	0.00257845639207701\\
-0.621621621621622	0.0025801795319527\\
-0.618618618618619	0.00258189473256846\\
-0.615615615615616	0.0025836020322782\\
-0.612612612612613	0.00258530146943583\\
-0.60960960960961	0.00258699308239528\\
-0.606606606606607	0.00258867690951044\\
-0.603603603603604	0.00259035298913523\\
-0.600600600600601	0.00259202135962356\\
-0.597597597597598	0.00259368205932936\\
-0.594594594594595	0.00259533512660652\\
-0.591591591591592	0.00259698059980896\\
-0.588588588588589	0.0025986185172906\\
-0.585585585585586	0.00260024891740535\\
-0.582582582582583	0.00260187183850712\\
-0.57957957957958	0.00260348731894982\\
-0.576576576576577	0.00260509539708737\\
-0.573573573573574	0.00260669611127368\\
-0.570570570570571	0.00260828949986266\\
-0.567567567567568	0.00260987560120822\\
-0.564564564564564	0.00261145445366428\\
-0.561561561561562	0.00261302609558475\\
-0.558558558558559	0.00261459056532354\\
-0.555555555555556	0.00261614790123457\\
-0.552552552552553	0.00261769814167174\\
-0.549549549549549	0.00261924132498898\\
-0.546546546546547	0.00262077748954019\\
-0.543543543543544	0.00262230667367928\\
-0.540540540540541	0.00262382891576017\\
-0.537537537537538	0.00262534425413678\\
-0.534534534534534	0.002626852727163\\
-0.531531531531532	0.00262835437319277\\
-0.528528528528529	0.00262984923057998\\
-0.525525525525526	0.00263133733767856\\
-0.522522522522523	0.00263281873284241\\
-0.519519519519519	0.00263429345442545\\
-0.516516516516517	0.00263576154078159\\
-0.513513513513513	0.00263722303026474\\
-0.510510510510511	0.00263867796122882\\
-0.507507507507508	0.00264012637202774\\
-0.504504504504504	0.00264156830101541\\
-0.501501501501502	0.00264300378654574\\
-0.498498498498498	0.00264443286697265\\
-0.495495495495496	0.00264585558065004\\
-0.492492492492492	0.00264727196593184\\
-0.489489489489489	0.00264868206117196\\
-0.486486486486487	0.0026500859047243\\
-0.483483483483483	0.00265148353494278\\
-0.480480480480481	0.00265287499018132\\
-0.477477477477477	0.00265426030879382\\
-0.474474474474474	0.0026556395291342\\
-0.471471471471472	0.00265701268955637\\
-0.468468468468468	0.00265837982841424\\
-0.465465465465465	0.00265974098406173\\
-0.462462462462462	0.00266109619485275\\
-0.459459459459459	0.00266244549914122\\
-0.456456456456456	0.00266378893528103\\
-0.453453453453453	0.00266512654162612\\
-0.45045045045045	0.00266645835653038\\
-0.447447447447447	0.00266778441834774\\
-0.444444444444444	0.0026691047654321\\
-0.441441441441441	0.00267041943613738\\
-0.438438438438438	0.00267172846881749\\
-0.435435435435435	0.00267303190182635\\
-0.432432432432432	0.00267432977351786\\
-0.429429429429429	0.00267562212224594\\
-0.426426426426426	0.0026769089863645\\
-0.423423423423423	0.00267819040422746\\
-0.42042042042042	0.00267946641418872\\
-0.417417417417417	0.0026807370546022\\
-0.414414414414414	0.00268200236382182\\
-0.411411411411411	0.00268326238020148\\
-0.408408408408408	0.0026845171420951\\
-0.405405405405405	0.00268576668785659\\
-0.402402402402402	0.00268701105583987\\
-0.399399399399399	0.00268825028439884\\
-0.396396396396396	0.00268948441188742\\
-0.393393393393393	0.00269071347665952\\
-0.39039039039039	0.00269193751706905\\
-0.387387387387387	0.00269315657146994\\
-0.384384384384384	0.00269437067821608\\
-0.381381381381381	0.00269557987566139\\
-0.378378378378378	0.00269678420215979\\
-0.375375375375375	0.00269798369606519\\
-0.372372372372372	0.0026991783957315\\
-0.369369369369369	0.00270036833951263\\
-0.366366366366366	0.0027015535657625\\
-0.363363363363363	0.00270273411283502\\
-0.36036036036036	0.00270391001908409\\
-0.357357357357357	0.00270508132286365\\
-0.354354354354354	0.00270624806252759\\
-0.351351351351351	0.00270741027642983\\
-0.348348348348348	0.00270856800292428\\
-0.345345345345345	0.00270972128036485\\
-0.342342342342342	0.00271087014710547\\
-0.339339339339339	0.00271201464150003\\
-0.336336336336336	0.00271315480190246\\
-0.333333333333333	0.00271429066666667\\
-0.33033033033033	0.00271542227414656\\
-0.327327327327327	0.00271654966269605\\
-0.324324324324324	0.00271767287066906\\
-0.321321321321321	0.0027187919364195\\
-0.318318318318318	0.00271990689830127\\
-0.315315315315315	0.0027210177946683\\
-0.312312312312312	0.00272212466387449\\
-0.309309309309309	0.00272322754427376\\
-0.306306306306306	0.00272432647422002\\
-0.303303303303303	0.00272542149206718\\
-0.3003003003003	0.00272651263616916\\
-0.297297297297297	0.00272759994487987\\
-0.294294294294294	0.00272868345655322\\
-0.291291291291291	0.00272976320954312\\
-0.288288288288288	0.00273083924220349\\
-0.285285285285285	0.00273191159288824\\
-0.282282282282282	0.00273298029995128\\
-0.279279279279279	0.00273404540174652\\
-0.276276276276276	0.00273510693662789\\
-0.273273273273273	0.00273616494294928\\
-0.27027027027027	0.00273721945906462\\
-0.267267267267267	0.00273827052332781\\
-0.264264264264264	0.00273931817409277\\
-0.261261261261261	0.00274036244971341\\
-0.258258258258258	0.00274140338854364\\
-0.255255255255255	0.00274244102893739\\
-0.252252252252252	0.00274347540924855\\
-0.249249249249249	0.00274450656783104\\
-0.246246246246246	0.00274553454303878\\
-0.243243243243243	0.00274655937322567\\
-0.24024024024024	0.00274758109674564\\
-0.237237237237237	0.00274859975195259\\
-0.234234234234234	0.00274961537720043\\
-0.231231231231231	0.00275062801084308\\
-0.228228228228228	0.00275163769123445\\
-0.225225225225225	0.00275264445672846\\
-0.222222222222222	0.00275364834567901\\
-0.219219219219219	0.00275464939644002\\
-0.216216216216216	0.00275564764736541\\
-0.213213213213213	0.00275664313680908\\
-0.21021021021021	0.00275763590312494\\
-0.207207207207207	0.00275862598466692\\
-0.204204204204204	0.00275961341978891\\
-0.201201201201201	0.00276059824684485\\
-0.198198198198198	0.00276158050418863\\
-0.195195195195195	0.00276256023017417\\
-0.192192192192192	0.00276353746315539\\
-0.189189189189189	0.00276451224148619\\
-0.186186186186186	0.00276548460352049\\
-0.183183183183183	0.0027664545876122\\
-0.18018018018018	0.00276742223211524\\
-0.177177177177177	0.00276838757538352\\
-0.174174174174174	0.00276935065577094\\
-0.171171171171171	0.00277031151163143\\
-0.168168168168168	0.00277127018131889\\
-0.165165165165165	0.00277222670318724\\
-0.162162162162162	0.00277318111559039\\
-0.159159159159159	0.00277413345688225\\
-0.156156156156156	0.00277508376541674\\
-0.153153153153153	0.00277603207954777\\
-0.15015015015015	0.00277697843762925\\
-0.147147147147147	0.0027779228780151\\
-0.144144144144144	0.00277886543905922\\
-0.141141141141141	0.00277980615911553\\
-0.138138138138138	0.00278074507653795\\
-0.135135135135135	0.00278168222968037\\
-0.132132132132132	0.00278261765689673\\
-0.129129129129129	0.00278355139654093\\
-0.126126126126126	0.00278448348696688\\
-0.123123123123123	0.0027854139665285\\
-0.12012012012012	0.00278634287357969\\
-0.117117117117117	0.00278727024647438\\
-0.114114114114114	0.00278819612356647\\
-0.111111111111111	0.00278912054320988\\
-0.108108108108108	0.00279004354375851\\
-0.105105105105105	0.00279096516356629\\
-0.102102102102102	0.00279188544098713\\
-0.0990990990990991	0.00279280441437493\\
-0.0960960960960962	0.00279372212208361\\
-0.0930930930930931	0.00279463860246709\\
-0.0900900900900901	0.00279555389387927\\
-0.0870870870870871	0.00279646803467407\\
-0.0840840840840841	0.0027973810632054\\
-0.081081081081081	0.00279829301782718\\
-0.0780780780780781	0.00279920393689331\\
-0.075075075075075	0.00280011385875771\\
-0.0720720720720721	0.00280102282177429\\
-0.0690690690690691	0.00280193086429697\\
-0.066066066066066	0.00280283802467966\\
-0.0630630630630631	0.00280374434127627\\
-0.06006006006006	0.0028046498524407\\
-0.0570570570570571	0.00280555459652689\\
-0.0540540540540541	0.00280645861188873\\
-0.051051051051051	0.00280736193688015\\
-0.0480480480480481	0.00280826460985505\\
-0.045045045045045	0.00280916666916734\\
-0.0420420420420421	0.00281006815317095\\
-0.039039039039039	0.00281096910021977\\
-0.036036036036036	0.00281186954866773\\
-0.0330330330330331	0.00281276953686874\\
-0.03003003003003	0.00281366910317671\\
-0.027027027027027	0.00281456828594555\\
-0.024024024024024	0.00281546712352918\\
-0.021021021021021	0.0028163656542815\\
-0.0180180180180181	0.00281726391655644\\
-0.015015015015015	0.0028181619487079\\
-0.012012012012012	0.0028190597890898\\
-0.00900900900900903	0.00281995747605604\\
-0.00600600600600598	0.00282085504796055\\
-0.00300300300300305	0.00282175254315723\\
0	0.00282265\\
0.00300300300300305	0.00282354745684277\\
0.00600600600600609	0.00282444495203945\\
0.00900900900900892	0.00282534252394396\\
0.012012012012012	0.0028262402109102\\
0.015015015015015	0.0028271380512921\\
0.0180180180180181	0.00282803608344356\\
0.0210210210210211	0.0028289343457185\\
0.0240240240240239	0.00282983287647082\\
0.027027027027027	0.00283073171405445\\
0.03003003003003	0.00283163089682329\\
0.0330330330330331	0.00283253046313126\\
0.0360360360360361	0.00283343045133227\\
0.0390390390390389	0.00283433089978023\\
0.042042042042042	0.00283523184682905\\
0.045045045045045	0.00283613333083266\\
0.0480480480480481	0.00283703539014495\\
0.0510510510510511	0.00283793806311985\\
0.0540540540540539	0.00283884138811127\\
0.057057057057057	0.00283974540347311\\
0.06006006006006	0.0028406501475593\\
0.0630630630630631	0.00284155565872373\\
0.0660660660660661	0.00284246197532034\\
0.0690690690690692	0.00284336913570303\\
0.072072072072072	0.00284427717822571\\
0.075075075075075	0.00284518614124229\\
0.0780780780780781	0.00284609606310669\\
0.0810810810810811	0.00284700698217282\\
0.0840840840840842	0.0028479189367946\\
0.087087087087087	0.00284883196532593\\
0.0900900900900901	0.00284974610612073\\
0.0930930930930931	0.00285066139753291\\
0.0960960960960962	0.00285157787791639\\
0.0990990990990992	0.00285249558562507\\
0.102102102102102	0.00285341455901287\\
0.105105105105105	0.00285433483643371\\
0.108108108108108	0.00285525645624149\\
0.111111111111111	0.00285617945679012\\
0.114114114114114	0.00285710387643353\\
0.117117117117117	0.00285802975352562\\
0.12012012012012	0.00285895712642031\\
0.123123123123123	0.0028598860334715\\
0.126126126126126	0.00286081651303312\\
0.129129129129129	0.00286174860345907\\
0.132132132132132	0.00286268234310327\\
0.135135135135135	0.00286361777031963\\
0.138138138138138	0.00286455492346205\\
0.141141141141141	0.00286549384088447\\
0.144144144144144	0.00286643456094078\\
0.147147147147147	0.0028673771219849\\
0.15015015015015	0.00286832156237075\\
0.153153153153153	0.00286926792045223\\
0.156156156156156	0.00287021623458326\\
0.159159159159159	0.00287116654311775\\
0.162162162162162	0.00287211888440961\\
0.165165165165165	0.00287307329681276\\
0.168168168168168	0.00287402981868111\\
0.171171171171171	0.00287498848836857\\
0.174174174174174	0.00287594934422906\\
0.177177177177177	0.00287691242461648\\
0.18018018018018	0.00287787776788476\\
0.183183183183183	0.0028788454123878\\
0.186186186186186	0.00287981539647951\\
0.189189189189189	0.00288078775851381\\
0.192192192192192	0.00288176253684461\\
0.195195195195195	0.00288273976982583\\
0.198198198198198	0.00288371949581137\\
0.201201201201201	0.00288470175315515\\
0.204204204204204	0.00288568658021109\\
0.207207207207207	0.00288667401533308\\
0.21021021021021	0.00288766409687506\\
0.213213213213213	0.00288865686319092\\
0.216216216216216	0.00288965235263459\\
0.219219219219219	0.00289065060355998\\
0.222222222222222	0.00289165165432099\\
0.225225225225225	0.00289265554327154\\
0.228228228228228	0.00289366230876555\\
0.231231231231231	0.00289467198915692\\
0.234234234234234	0.00289568462279957\\
0.237237237237237	0.00289670024804741\\
0.24024024024024	0.00289771890325436\\
0.243243243243243	0.00289874062677433\\
0.246246246246246	0.00289976545696122\\
0.249249249249249	0.00290079343216896\\
0.252252252252252	0.00290182459075145\\
0.255255255255255	0.00290285897106261\\
0.258258258258258	0.00290389661145636\\
0.261261261261261	0.00290493755028659\\
0.264264264264264	0.00290598182590723\\
0.267267267267267	0.00290702947667219\\
0.27027027027027	0.00290808054093538\\
0.273273273273273	0.00290913505705072\\
0.276276276276276	0.00291019306337211\\
0.279279279279279	0.00291125459825348\\
0.282282282282282	0.00291231970004872\\
0.285285285285285	0.00291338840711176\\
0.288288288288288	0.00291446075779651\\
0.291291291291291	0.00291553679045688\\
0.294294294294294	0.00291661654344678\\
0.297297297297297	0.00291770005512013\\
0.3003003003003	0.00291878736383084\\
0.303303303303303	0.00291987850793282\\
0.306306306306306	0.00292097352577998\\
0.309309309309309	0.00292207245572624\\
0.312312312312312	0.00292317533612551\\
0.315315315315315	0.0029242822053317\\
0.318318318318318	0.00292539310169873\\
0.321321321321321	0.0029265080635805\\
0.324324324324324	0.00292762712933094\\
0.327327327327327	0.00292875033730395\\
0.33033033033033	0.00292987772585344\\
0.333333333333333	0.00293100933333333\\
0.336336336336336	0.00293214519809754\\
0.339339339339339	0.00293328535849997\\
0.342342342342342	0.00293442985289453\\
0.345345345345345	0.00293557871963515\\
0.348348348348348	0.00293673199707572\\
0.351351351351351	0.00293788972357017\\
0.354354354354354	0.00293905193747241\\
0.357357357357357	0.00294021867713635\\
0.36036036036036	0.0029413899809159\\
0.363363363363363	0.00294256588716498\\
0.366366366366366	0.0029437464342375\\
0.369369369369369	0.00294493166048737\\
0.372372372372372	0.0029461216042685\\
0.375375375375375	0.00294731630393481\\
0.378378378378378	0.00294851579784021\\
0.381381381381381	0.00294972012433861\\
0.384384384384384	0.00295092932178392\\
0.387387387387387	0.00295214342853006\\
0.39039039039039	0.00295336248293095\\
0.393393393393393	0.00295458652334048\\
0.396396396396396	0.00295581558811258\\
0.399399399399399	0.00295704971560116\\
0.402402402402402	0.00295828894416013\\
0.405405405405405	0.00295953331214341\\
0.408408408408408	0.0029607828579049\\
0.411411411411411	0.00296203761979852\\
0.414414414414414	0.00296329763617818\\
0.417417417417417	0.0029645629453978\\
0.42042042042042	0.00296583358581128\\
0.423423423423423	0.00296710959577254\\
0.426426426426426	0.0029683910136355\\
0.429429429429429	0.00296967787775406\\
0.432432432432432	0.00297097022648214\\
0.435435435435436	0.00297226809817365\\
0.438438438438439	0.00297357153118251\\
0.441441441441441	0.00297488056386262\\
0.444444444444444	0.0029761952345679\\
0.447447447447447	0.00297751558165226\\
0.450450450450451	0.00297884164346962\\
0.453453453453454	0.00298017345837388\\
0.456456456456456	0.00298151106471897\\
0.459459459459459	0.00298285450085878\\
0.462462462462462	0.00298420380514725\\
0.465465465465466	0.00298555901593827\\
0.468468468468469	0.00298692017158576\\
0.471471471471471	0.00298828731044363\\
0.474474474474474	0.0029896604708658\\
0.477477477477477	0.00299103969120618\\
0.480480480480481	0.00299242500981868\\
0.483483483483484	0.00299381646505722\\
0.486486486486486	0.0029952140952757\\
0.489489489489489	0.00299661793882804\\
0.492492492492492	0.00299802803406816\\
0.495495495495496	0.00299944441934996\\
0.498498498498499	0.00300086713302735\\
0.501501501501501	0.00300229621345426\\
0.504504504504504	0.00300373169898459\\
0.507507507507508	0.00300517362797226\\
0.510510510510511	0.00300662203877118\\
0.513513513513514	0.00300807696973526\\
0.516516516516516	0.00300953845921841\\
0.519519519519519	0.00301100654557455\\
0.522522522522523	0.00301248126715759\\
0.525525525525526	0.00301396266232144\\
0.528528528528529	0.00301545076942002\\
0.531531531531531	0.00301694562680723\\
0.534534534534534	0.003018447272837\\
0.537537537537538	0.00301995574586322\\
0.540540540540541	0.00302147108423983\\
0.543543543543544	0.00302299332632072\\
0.546546546546546	0.00302452251045981\\
0.549549549549549	0.00302605867501102\\
0.552552552552553	0.00302760185832826\\
0.555555555555556	0.00302915209876543\\
0.558558558558559	0.00303070943467646\\
0.561561561561561	0.00303227390441525\\
0.564564564564564	0.00303384554633572\\
0.567567567567568	0.00303542439879178\\
0.570570570570571	0.00303701050013734\\
0.573573573573574	0.00303860388872632\\
0.576576576576577	0.00304020460291263\\
0.57957957957958	0.00304181268105018\\
0.582582582582583	0.00304342816149288\\
0.585585585585586	0.00304505108259465\\
0.588588588588589	0.0030466814827094\\
0.591591591591592	0.00304831940019104\\
0.594594594594595	0.00304996487339348\\
0.597597597597598	0.00305161794067064\\
0.600600600600601	0.00305327864037644\\
0.603603603603604	0.00305494701086477\\
0.606606606606607	0.00305662309048956\\
0.60960960960961	0.00305830691760472\\
0.612612612612613	0.00305999853056417\\
0.615615615615616	0.0030616979677218\\
0.618618618618619	0.00306340526743154\\
0.621621621621622	0.0030651204680473\\
0.624624624624625	0.00306684360792299\\
0.627627627627628	0.00306857472541253\\
0.630630630630631	0.00307031385886983\\
0.633633633633634	0.00307206104664879\\
0.636636636636637	0.00307381632710334\\
0.63963963963964	0.00307557973858738\\
0.642642642642643	0.00307735131945483\\
0.645645645645646	0.00307913110805961\\
0.648648648648649	0.00308091914275561\\
0.651651651651652	0.00308271546189676\\
0.654654654654655	0.00308452010383698\\
0.657657657657658	0.00308633310693016\\
0.660660660660661	0.00308815450953023\\
0.663663663663664	0.00308998434999109\\
0.666666666666667	0.00309182266666667\\
0.66966966966967	0.00309366949791086\\
0.672672672672673	0.0030955248820776\\
0.675675675675676	0.00309738885752078\\
0.678678678678679	0.00309926146259432\\
0.681681681681682	0.00310114273565214\\
0.684684684684685	0.00310303271504814\\
0.687687687687688	0.00310493143913624\\
0.690690690690691	0.00310683894627035\\
0.693693693693694	0.00310875527480439\\
0.696696696696697	0.00311068046309226\\
0.6996996996997	0.00311261454948789\\
0.702702702702703	0.00311455757234517\\
0.705705705705706	0.00311650957001803\\
0.708708708708709	0.00311847058086038\\
0.711711711711712	0.00312044064322613\\
0.714714714714715	0.0031224197954692\\
0.717717717717718	0.00312440807594349\\
0.720720720720721	0.00312640552300291\\
0.723723723723724	0.0031284121750014\\
0.726726726726727	0.00313042807029284\\
0.72972972972973	0.00313245324723116\\
0.732732732732733	0.00313448774417027\\
0.735735735735736	0.00313653159946409\\
0.738738738738739	0.00313858485146651\\
0.741741741741742	0.00314064753853147\\
0.744744744744745	0.00314271969901287\\
0.747747747747748	0.00314480137126462\\
0.750750750750751	0.00314689259364063\\
0.753753753753754	0.00314899340449483\\
0.756756756756757	0.00315110384218111\\
0.75975975975976	0.00315322394505341\\
0.762762762762763	0.00315535375146561\\
0.765765765765766	0.00315749329977165\\
0.768768768768769	0.00315964262832543\\
0.771771771771772	0.00316180177548086\\
0.774774774774775	0.00316397077959186\\
0.777777777777778	0.00316614967901235\\
0.780780780780781	0.00316833851209622\\
0.783783783783784	0.0031705373171974\\
0.786786786786787	0.0031727461326698\\
0.78978978978979	0.00317496499686733\\
0.792792792792793	0.00317719394814391\\
0.795795795795796	0.00317943302485344\\
0.798798798798799	0.00318168226534984\\
0.801801801801802	0.00318394170798702\\
0.804804804804805	0.0031862113911189\\
0.807807807807808	0.00318849135309939\\
0.810810810810811	0.00319078163228239\\
0.813813813813814	0.00319308226702183\\
0.816816816816817	0.00319539329567162\\
0.81981981981982	0.00319771475658566\\
0.822822822822823	0.00320004668811787\\
0.825825825825826	0.00320238912862217\\
0.828828828828829	0.00320474211645246\\
0.831831831831832	0.00320710568996267\\
0.834834834834835	0.00320947988750669\\
0.837837837837838	0.00321186474743845\\
0.840840840840841	0.00321426030811186\\
0.843843843843844	0.00321666660788083\\
0.846846846846847	0.00321908368509927\\
0.84984984984985	0.0032215115781211\\
0.852852852852853	0.00322395032530022\\
0.855855855855856	0.00322639996499056\\
0.858858858858859	0.00322886053554602\\
0.861861861861862	0.00323133207532051\\
0.864864864864865	0.00323381462266796\\
0.867867867867868	0.00323630821594226\\
0.870870870870871	0.00323881289349734\\
0.873873873873874	0.00324132869368711\\
0.876876876876877	0.00324385565486548\\
0.87987987987988	0.00324639381538636\\
0.882882882882883	0.00324894321360367\\
0.885885885885886	0.00325150388787131\\
0.888888888888889	0.00325407587654321\\
0.891891891891892	0.00325665921797327\\
0.894894894894895	0.0032592539505154\\
0.897897897897898	0.00326186011252353\\
0.900900900900901	0.00326447774235155\\
0.903903903903904	0.0032671068783534\\
0.906906906906907	0.00326974755888296\\
0.90990990990991	0.00327239982229417\\
0.912912912912913	0.00327506370694093\\
0.915915915915916	0.00327773925117716\\
0.918918918918919	0.00328042649335676\\
0.921921921921922	0.00328312547183365\\
0.924924924924925	0.00328583622496175\\
0.927927927927928	0.00328855879109497\\
0.930930930930931	0.00329129320858721\\
0.933933933933934	0.00329403951579239\\
0.936936936936937	0.00329679775106443\\
0.93993993993994	0.00329956795275724\\
0.942942942942943	0.00330235015922472\\
0.945945945945946	0.0033051444088208\\
0.948948948948949	0.00330795073989938\\
0.951951951951952	0.00331076919081438\\
0.954954954954955	0.00331359979991972\\
0.957957957957958	0.00331644260556929\\
0.960960960960961	0.00331929764611702\\
0.963963963963964	0.00332216495991682\\
0.966966966966967	0.0033250445853226\\
0.96996996996997	0.00332793656068827\\
0.972972972972973	0.00333084092436776\\
0.975975975975976	0.00333375771471496\\
0.978978978978979	0.00333668697008379\\
0.981981981981982	0.00333962872882817\\
0.984984984984985	0.00334258302930201\\
0.987987987987988	0.00334554990985921\\
0.990990990990991	0.0033485294088537\\
0.993993993993994	0.00335152156463939\\
0.996996996996997	0.00335452641557018\\
1	0.003357544\\
1.003003003003	0.00336057435628275\\
1.00600600600601	0.00336361752277235\\
1.00900900900901	0.0033666735378227\\
1.01201201201201	0.00336974243978773\\
1.01501501501502	0.00337282426702134\\
1.01801801801802	0.00337591905787745\\
1.02102102102102	0.00337902685070997\\
1.02402402402402	0.00338214768387282\\
1.02702702702703	0.0033852815957199\\
1.03003003003003	0.00338842862460513\\
1.03303303303303	0.00339158880888242\\
1.03603603603604	0.00339476218690568\\
1.03903903903904	0.00339794879702883\\
1.04204204204204	0.00340114867760578\\
1.04504504504505	0.00340436186699044\\
1.04804804804805	0.00340758840353672\\
1.05105105105105	0.00341082832559855\\
1.05405405405405	0.00341408167152982\\
1.05705705705706	0.00341734847968446\\
1.06006006006006	0.00342062878841637\\
1.06306306306306	0.00342392263607947\\
1.06606606606607	0.00342723006102767\\
1.06906906906907	0.00343055110161488\\
1.07207207207207	0.00343388579619503\\
1.07507507507508	0.00343723418312201\\
1.07807807807808	0.00344059630074974\\
1.08108108108108	0.00344397218743214\\
1.08408408408408	0.00344736188152311\\
1.08708708708709	0.00345076542137657\\
1.09009009009009	0.00345418284534644\\
1.09309309309309	0.00345761419178662\\
1.0960960960961	0.00346105949905104\\
1.0990990990991	0.00346451880549359\\
1.1021021021021	0.00346799214946819\\
1.10510510510511	0.00347147956932876\\
1.10810810810811	0.00347498110342921\\
1.11111111111111	0.00347849679012346\\
1.11411411411411	0.0034820266677654\\
1.11711711711712	0.00348557077470897\\
1.12012012012012	0.00348912914930806\\
1.12312312312312	0.0034927018299166\\
1.12612612612613	0.00349628885488849\\
1.12912912912913	0.00349989026257765\\
1.13213213213213	0.00350350609133799\\
1.13513513513514	0.00350713637952342\\
1.13813813813814	0.00351078116548786\\
1.14114114114114	0.00351444048758522\\
1.14414414414414	0.00351811438416941\\
1.14714714714715	0.00352180289359435\\
1.15015015015015	0.00352550605421394\\
1.15315315315315	0.0035292239043821\\
1.15615615615616	0.00353295648245275\\
1.15915915915916	0.00353670382677979\\
1.16216216216216	0.00354046597571713\\
1.16516516516517	0.0035442429676187\\
1.16816816816817	0.0035480348408384\\
1.17117117117117	0.00355184163373015\\
1.17417417417417	0.00355566338464786\\
1.17717717717718	0.00355950013194543\\
1.18018018018018	0.0035633519139768\\
1.18318318318318	0.00356721876909586\\
1.18618618618619	0.00357110073565652\\
1.18918918918919	0.00357499785201271\\
1.19219219219219	0.00357891015651834\\
1.1951951951952	0.00358283768752731\\
1.1981981981982	0.00358678048339355\\
1.2012012012012	0.00359073858247095\\
1.2042042042042	0.00359471202311345\\
1.20720720720721	0.00359870084367494\\
1.21021021021021	0.00360270508250934\\
1.21321321321321	0.00360672477797057\\
1.21621621621622	0.00361075996841253\\
1.21921921921922	0.00361481069218915\\
1.22222222222222	0.00361887698765432\\
1.22522522522523	0.00362295889316197\\
1.22822822822823	0.00362705644706601\\
1.23123123123123	0.00363116968772035\\
1.23423423423423	0.0036352986534789\\
1.23723723723724	0.00363944338269558\\
1.24024024024024	0.00364360391372429\\
1.24324324324324	0.00364778028491896\\
1.24624624624625	0.00365197253463349\\
1.24924924924925	0.0036561807012218\\
1.25225225225225	0.00366040482303779\\
1.25525525525526	0.00366464493843539\\
1.25825825825826	0.00366890108576851\\
1.26126126126126	0.00367317330339105\\
1.26426426426426	0.00367746162965693\\
1.26726726726727	0.00368176610292006\\
1.27027027027027	0.00368608676153436\\
1.27327327327327	0.00369042364385374\\
1.27627627627628	0.00369477678823211\\
1.27927927927928	0.00369914623302338\\
1.28228228228228	0.00370353201658147\\
1.28528528528529	0.00370793417726029\\
1.28828828828829	0.00371235275341375\\
1.29129129129129	0.00371678778339576\\
1.29429429429429	0.00372123930556025\\
1.2972972972973	0.00372570735826111\\
1.3003003003003	0.00373019197985227\\
1.3033033033033	0.00373469320868763\\
1.30630630630631	0.00373921108312111\\
1.30930930930931	0.00374374564150661\\
1.31231231231231	0.00374829692219807\\
1.31531531531532	0.00375286496354938\\
1.31831831831832	0.00375744980391446\\
1.32132132132132	0.00376205148164722\\
1.32432432432432	0.00376667003510157\\
1.32732732732733	0.00377130550263144\\
1.33033033033033	0.00377595792259072\\
1.33333333333333	0.00378062733333333\\
1.33633633633634	0.00378531377321319\\
1.33933933933934	0.00379001728058421\\
1.34234234234234	0.0037947378938003\\
1.34534534534535	0.00379947565121537\\
1.34834834834835	0.00380423059118334\\
1.35135135135135	0.00380900275205812\\
1.35435435435435	0.00381379217219362\\
1.35735735735736	0.00381859888994375\\
1.36036036036036	0.00382342294366244\\
1.36336336336336	0.00382826437170358\\
1.36636636636637	0.00383312321242109\\
1.36936936936937	0.00383799950416889\\
1.37237237237237	0.00384289328530088\\
1.37537537537538	0.00384780459417099\\
1.37837837837838	0.00385273346913312\\
1.38138138138138	0.00385767994854119\\
1.38438438438438	0.0038626440707491\\
1.38738738738739	0.00386762587411078\\
1.39039039039039	0.00387262539698013\\
1.39339339339339	0.00387764267771107\\
1.3963963963964	0.00388267775465751\\
1.3993993993994	0.00388773066617336\\
1.4024024024024	0.00389280145061253\\
1.40540540540541	0.00389789014632894\\
1.40840840840841	0.00390299679167651\\
1.41141141141141	0.00390812142500913\\
1.41441441441441	0.00391326408468074\\
1.41741741741742	0.00391842480904523\\
1.42042042042042	0.00392360363645652\\
1.42342342342342	0.00392880060526853\\
1.42642642642643	0.00393401575383516\\
1.42942942942943	0.00393924912051033\\
1.43243243243243	0.00394450074364796\\
1.43543543543544	0.00394977066160195\\
1.43843843843844	0.00395505891272621\\
1.44144144144144	0.00396036553537467\\
1.44444444444444	0.00396569056790123\\
1.44744744744745	0.00397103404865981\\
1.45045045045045	0.00397639601600432\\
1.45345345345345	0.00398177650828866\\
1.45645645645646	0.00398717556386676\\
1.45945945945946	0.00399259322109253\\
1.46246246246246	0.00399802951831988\\
1.46546546546547	0.00400348449390272\\
1.46846846846847	0.00400895818619496\\
1.47147147147147	0.00401445063355052\\
1.47447447447447	0.00401996187432331\\
1.47747747747748	0.00402549194686725\\
1.48048048048048	0.00403104088953624\\
1.48348348348348	0.0040366087406842\\
1.48648648648649	0.00404219553866503\\
1.48948948948949	0.00404780132183267\\
1.49249249249249	0.00405342612854101\\
1.4954954954955	0.00405906999714397\\
1.4984984984985	0.00406473296599546\\
1.5015015015015	0.00407041507344939\\
1.5045045045045	0.00407611635785969\\
1.50750750750751	0.00408183685758025\\
1.51051051051051	0.00408757661096499\\
1.51351351351351	0.00409333565636784\\
1.51651651651652	0.00409911403214269\\
1.51951951951952	0.00410491177664346\\
1.52252252252252	0.00411072892822406\\
1.52552552552553	0.00411656552523841\\
1.52852852852853	0.00412242160604042\\
1.53153153153153	0.004128297208984\\
1.53453453453453	0.00413419237242307\\
1.53753753753754	0.00414010713471153\\
1.54054054054054	0.0041460415342033\\
1.54354354354354	0.0041519956092523\\
1.54654654654655	0.00415796939821243\\
1.54954954954955	0.00416396293943761\\
1.55255255255255	0.00416997627128175\\
1.55555555555556	0.00417600943209876\\
1.55855855855856	0.00418206246024257\\
1.56156156156156	0.00418813539406706\\
1.56456456456456	0.00419422827192618\\
1.56756756756757	0.00420034113217381\\
1.57057057057057	0.00420647401316388\\
1.57357357357357	0.0042126269532503\\
1.57657657657658	0.00421879999078699\\
1.57957957957958	0.00422499316412785\\
1.58258258258258	0.0042312065116268\\
1.58558558558559	0.00423744007163774\\
1.58858858858859	0.00424369388251461\\
1.59159159159159	0.00424996798261129\\
1.59459459459459	0.00425626241028172\\
1.5975975975976	0.0042625772038798\\
1.6006006006006	0.00426891240175944\\
1.6036036036036	0.00427526804227456\\
1.60660660660661	0.00428164416377907\\
1.60960960960961	0.00428804080462688\\
1.61261261261261	0.00429445800317191\\
1.61561561561562	0.00430089579776806\\
1.61861861861862	0.00430735422676926\\
1.62162162162162	0.00431383332852941\\
1.62462462462462	0.00432033314140242\\
1.62762762762763	0.00432685370374221\\
1.63063063063063	0.0043333950539027\\
1.63363363363363	0.00433995723023779\\
1.63663663663664	0.00434654027110139\\
1.63963963963964	0.00435314421484743\\
1.64264264264264	0.0043597690998298\\
1.64564564564565	0.00436641496440244\\
1.64864864864865	0.00437308184691923\\
1.65165165165165	0.00437976978573411\\
1.65465465465465	0.00438647881920099\\
1.65765765765766	0.00439320898567377\\
1.66066066066066	0.00439996032350637\\
1.66366366366366	0.00440673287105269\\
1.66666666666667	0.00441352666666667\\
1.66966966966967	0.0044203417487022\\
1.67267267267267	0.00442717815551319\\
1.67567567567568	0.00443403592545358\\
1.67867867867868	0.00444091509687725\\
1.68168168168168	0.00444781570813814\\
1.68468468468468	0.00445473779759014\\
1.68768768768769	0.00446168140358718\\
1.69069069069069	0.00446864656448316\\
1.69369369369369	0.004475633318632\\
1.6966966966967	0.00448264170438761\\
1.6996996996997	0.00448967176010391\\
1.7027027027027	0.0044967235241348\\
1.70570570570571	0.0045037970348342\\
1.70870870870871	0.00451089233055602\\
1.71171171171171	0.00451800944965418\\
1.71471471471471	0.00452514843048259\\
1.71771771771772	0.00453230931139515\\
1.72072072072072	0.00453949213074579\\
1.72372372372372	0.00454669692688842\\
1.72672672672673	0.00455392373817694\\
1.72972972972973	0.00456117260296527\\
1.73273273273273	0.00456844355960733\\
1.73573573573574	0.00457573664645703\\
1.73873873873874	0.00458305190186827\\
1.74174174174174	0.00459038936419497\\
1.74474474474474	0.00459774907179105\\
1.74774774774775	0.00460513106301042\\
1.75075075075075	0.00461253537620698\\
1.75375375375375	0.00461996204973466\\
1.75675675675676	0.00462741112194737\\
1.75975975975976	0.00463488263119901\\
1.76276276276276	0.0046423766158435\\
1.76576576576577	0.00464989311423476\\
1.76876876876877	0.00465743216472669\\
1.77177177177177	0.00466499380567322\\
1.77477477477477	0.00467257807542824\\
1.77777777777778	0.00468018501234568\\
1.78078078078078	0.00468781465477944\\
1.78378378378378	0.00469546704108345\\
1.78678678678679	0.00470314220961161\\
1.78978978978979	0.00471084019871783\\
1.79279279279279	0.00471856104675603\\
1.7957957957958	0.00472630479208013\\
1.7987987987988	0.00473407147304402\\
1.8018018018018	0.00474186112800163\\
1.8048048048048	0.00474967379530687\\
1.80780780780781	0.00475750951331365\\
1.81081081081081	0.00476536832037589\\
1.81381381381381	0.00477325025484749\\
1.81681681681682	0.00478115535508238\\
1.81981981981982	0.00478908365943445\\
1.82282282282282	0.00479703520625763\\
1.82582582582583	0.00480501003390583\\
1.82882882882883	0.00481300818073296\\
1.83183183183183	0.00482102968509293\\
1.83483483483483	0.00482907458533966\\
1.83783783783784	0.00483714291982706\\
1.84084084084084	0.00484523472690904\\
1.84384384384384	0.00485335004493951\\
1.84684684684685	0.00486148891227239\\
1.84984984984985	0.00486965136726159\\
1.85285285285285	0.00487783744826102\\
1.85585585585586	0.0048860471936246\\
1.85885885885886	0.00489428064170623\\
1.86186186186186	0.00490253783085983\\
1.86486486486486	0.00491081879943932\\
1.86786786786787	0.00491912358579861\\
1.87087087087087	0.0049274522282916\\
1.87387387387387	0.00493580476527221\\
1.87687687687688	0.00494418123509436\\
1.87987987987988	0.00495258167611195\\
1.88288288288288	0.00496100612667891\\
1.88588588588589	0.00496945462514913\\
1.88888888888889	0.00497792720987654\\
1.89189189189189	0.00498642391921505\\
1.89489489489489	0.00499494479151857\\
1.8978978978979	0.00500348986514101\\
1.9009009009009	0.00501205917843629\\
1.9039039039039	0.00502065276975831\\
1.90690690690691	0.005029270677461\\
1.90990990990991	0.00503791293989826\\
1.91291291291291	0.00504657959542401\\
1.91591591591592	0.00505527068239216\\
1.91891891891892	0.00506398623915661\\
1.92192192192192	0.0050727263040713\\
1.92492492492492	0.00508149091549012\\
1.92792792792793	0.00509028011176699\\
1.93093093093093	0.00509909393125582\\
1.93393393393393	0.00510793241231053\\
1.93693693693694	0.00511679559328503\\
1.93993993993994	0.00512568351253323\\
1.94294294294294	0.00513459620840904\\
1.94594594594595	0.00514353371926638\\
1.94894894894895	0.00515249608345916\\
1.95195195195195	0.00516148333934129\\
1.95495495495495	0.00517049552526668\\
1.95795795795796	0.00517953267958926\\
1.96096096096096	0.00518859484066292\\
1.96396396396396	0.00519768204684158\\
1.96696696696697	0.00520679433647916\\
1.96996996996997	0.00521593174792957\\
1.97297297297297	0.00522509431954672\\
1.97597597597598	0.00523428208968452\\
1.97897897897898	0.00524349509669689\\
1.98198198198198	0.00525273337893774\\
1.98498498498498	0.00526199697476098\\
1.98798798798799	0.00527128592252052\\
1.99099099099099	0.00528060026057029\\
1.99399399399399	0.00528994002726418\\
1.996996996997	0.00529930526095611\\
2	0.005308696\\
};
\addlegendentry{Fit}

\addplot [color=black, dashed]
  table[row sep=crcr]{%
-1	0.002208196\\
-0.996996996996997	0.00220987114714715\\
-0.993993993993994	0.00221154629429429\\
-0.990990990990991	0.00221322144144144\\
-0.987987987987988	0.00221489658858859\\
-0.984984984984985	0.00221657173573574\\
-0.981981981981982	0.00221824688288288\\
-0.978978978978979	0.00221992203003003\\
-0.975975975975976	0.00222159717717718\\
-0.972972972972973	0.00222327232432432\\
-0.96996996996997	0.00222494747147147\\
-0.966966966966967	0.00222662261861862\\
-0.963963963963964	0.00222829776576577\\
-0.960960960960961	0.00222997291291291\\
-0.957957957957958	0.00223164806006006\\
-0.954954954954955	0.00223332320720721\\
-0.951951951951952	0.00223499835435435\\
-0.948948948948949	0.0022366735015015\\
-0.945945945945946	0.00223834864864865\\
-0.942942942942943	0.0022400237957958\\
-0.93993993993994	0.00224169894294294\\
-0.936936936936937	0.00224337409009009\\
-0.933933933933934	0.00224504923723724\\
-0.930930930930931	0.00224672438438438\\
-0.927927927927928	0.00224839953153153\\
-0.924924924924925	0.00225007467867868\\
-0.921921921921922	0.00225174982582583\\
-0.918918918918919	0.00225342497297297\\
-0.915915915915916	0.00225510012012012\\
-0.912912912912913	0.00225677526726727\\
-0.90990990990991	0.00225845041441441\\
-0.906906906906907	0.00226012556156156\\
-0.903903903903904	0.00226180070870871\\
-0.900900900900901	0.00226347585585586\\
-0.897897897897898	0.002265151003003\\
-0.894894894894895	0.00226682615015015\\
-0.891891891891892	0.0022685012972973\\
-0.888888888888889	0.00227017644444444\\
-0.885885885885886	0.00227185159159159\\
-0.882882882882883	0.00227352673873874\\
-0.87987987987988	0.00227520188588589\\
-0.876876876876877	0.00227687703303303\\
-0.873873873873874	0.00227855218018018\\
-0.870870870870871	0.00228022732732733\\
-0.867867867867868	0.00228190247447447\\
-0.864864864864865	0.00228357762162162\\
-0.861861861861862	0.00228525276876877\\
-0.858858858858859	0.00228692791591592\\
-0.855855855855856	0.00228860306306306\\
-0.852852852852853	0.00229027821021021\\
-0.84984984984985	0.00229195335735736\\
-0.846846846846847	0.0022936285045045\\
-0.843843843843844	0.00229530365165165\\
-0.840840840840841	0.0022969787987988\\
-0.837837837837838	0.00229865394594595\\
-0.834834834834835	0.00230032909309309\\
-0.831831831831832	0.00230200424024024\\
-0.828828828828829	0.00230367938738739\\
-0.825825825825826	0.00230535453453453\\
-0.822822822822823	0.00230702968168168\\
-0.81981981981982	0.00230870482882883\\
-0.816816816816817	0.00231037997597598\\
-0.813813813813814	0.00231205512312312\\
-0.810810810810811	0.00231373027027027\\
-0.807807807807808	0.00231540541741742\\
-0.804804804804805	0.00231708056456456\\
-0.801801801801802	0.00231875571171171\\
-0.798798798798799	0.00232043085885886\\
-0.795795795795796	0.00232210600600601\\
-0.792792792792793	0.00232378115315315\\
-0.78978978978979	0.0023254563003003\\
-0.786786786786787	0.00232713144744745\\
-0.783783783783784	0.00232880659459459\\
-0.780780780780781	0.00233048174174174\\
-0.777777777777778	0.00233215688888889\\
-0.774774774774775	0.00233383203603604\\
-0.771771771771772	0.00233550718318318\\
-0.768768768768769	0.00233718233033033\\
-0.765765765765766	0.00233885747747748\\
-0.762762762762763	0.00234053262462462\\
-0.75975975975976	0.00234220777177177\\
-0.756756756756757	0.00234388291891892\\
-0.753753753753754	0.00234555806606607\\
-0.750750750750751	0.00234723321321321\\
-0.747747747747748	0.00234890836036036\\
-0.744744744744745	0.00235058350750751\\
-0.741741741741742	0.00235225865465465\\
-0.738738738738739	0.0023539338018018\\
-0.735735735735736	0.00235560894894895\\
-0.732732732732733	0.0023572840960961\\
-0.72972972972973	0.00235895924324324\\
-0.726726726726727	0.00236063439039039\\
-0.723723723723724	0.00236230953753754\\
-0.720720720720721	0.00236398468468468\\
-0.717717717717718	0.00236565983183183\\
-0.714714714714715	0.00236733497897898\\
-0.711711711711712	0.00236901012612613\\
-0.708708708708709	0.00237068527327327\\
-0.705705705705706	0.00237236042042042\\
-0.702702702702703	0.00237403556756757\\
-0.6996996996997	0.00237571071471471\\
-0.696696696696697	0.00237738586186186\\
-0.693693693693694	0.00237906100900901\\
-0.690690690690691	0.00238073615615616\\
-0.687687687687688	0.0023824113033033\\
-0.684684684684685	0.00238408645045045\\
-0.681681681681682	0.0023857615975976\\
-0.678678678678679	0.00238743674474474\\
-0.675675675675676	0.00238911189189189\\
-0.672672672672673	0.00239078703903904\\
-0.66966966966967	0.00239246218618619\\
-0.666666666666667	0.00239413733333333\\
-0.663663663663664	0.00239581248048048\\
-0.660660660660661	0.00239748762762763\\
-0.657657657657658	0.00239916277477477\\
-0.654654654654655	0.00240083792192192\\
-0.651651651651652	0.00240251306906907\\
-0.648648648648649	0.00240418821621622\\
-0.645645645645646	0.00240586336336336\\
-0.642642642642643	0.00240753851051051\\
-0.63963963963964	0.00240921365765766\\
-0.636636636636637	0.00241088880480481\\
-0.633633633633634	0.00241256395195195\\
-0.630630630630631	0.0024142390990991\\
-0.627627627627628	0.00241591424624625\\
-0.624624624624625	0.00241758939339339\\
-0.621621621621622	0.00241926454054054\\
-0.618618618618619	0.00242093968768769\\
-0.615615615615616	0.00242261483483483\\
-0.612612612612613	0.00242428998198198\\
-0.60960960960961	0.00242596512912913\\
-0.606606606606607	0.00242764027627628\\
-0.603603603603604	0.00242931542342342\\
-0.600600600600601	0.00243099057057057\\
-0.597597597597598	0.00243266571771772\\
-0.594594594594595	0.00243434086486487\\
-0.591591591591592	0.00243601601201201\\
-0.588588588588589	0.00243769115915916\\
-0.585585585585586	0.00243936630630631\\
-0.582582582582583	0.00244104145345345\\
-0.57957957957958	0.0024427166006006\\
-0.576576576576577	0.00244439174774775\\
-0.573573573573574	0.00244606689489489\\
-0.570570570570571	0.00244774204204204\\
-0.567567567567568	0.00244941718918919\\
-0.564564564564564	0.00245109233633634\\
-0.561561561561562	0.00245276748348348\\
-0.558558558558559	0.00245444263063063\\
-0.555555555555556	0.00245611777777778\\
-0.552552552552553	0.00245779292492493\\
-0.549549549549549	0.00245946807207207\\
-0.546546546546547	0.00246114321921922\\
-0.543543543543544	0.00246281836636637\\
-0.540540540540541	0.00246449351351351\\
-0.537537537537538	0.00246616866066066\\
-0.534534534534534	0.00246784380780781\\
-0.531531531531532	0.00246951895495495\\
-0.528528528528529	0.0024711941021021\\
-0.525525525525526	0.00247286924924925\\
-0.522522522522523	0.0024745443963964\\
-0.519519519519519	0.00247621954354354\\
-0.516516516516517	0.00247789469069069\\
-0.513513513513513	0.00247956983783784\\
-0.510510510510511	0.00248124498498499\\
-0.507507507507508	0.00248292013213213\\
-0.504504504504504	0.00248459527927928\\
-0.501501501501502	0.00248627042642643\\
-0.498498498498498	0.00248794557357357\\
-0.495495495495496	0.00248962072072072\\
-0.492492492492492	0.00249129586786787\\
-0.489489489489489	0.00249297101501502\\
-0.486486486486487	0.00249464616216216\\
-0.483483483483483	0.00249632130930931\\
-0.480480480480481	0.00249799645645646\\
-0.477477477477477	0.0024996716036036\\
-0.474474474474474	0.00250134675075075\\
-0.471471471471472	0.0025030218978979\\
-0.468468468468468	0.00250469704504505\\
-0.465465465465465	0.00250637219219219\\
-0.462462462462462	0.00250804733933934\\
-0.459459459459459	0.00250972248648649\\
-0.456456456456456	0.00251139763363363\\
-0.453453453453453	0.00251307278078078\\
-0.45045045045045	0.00251474792792793\\
-0.447447447447447	0.00251642307507508\\
-0.444444444444444	0.00251809822222222\\
-0.441441441441441	0.00251977336936937\\
-0.438438438438438	0.00252144851651652\\
-0.435435435435435	0.00252312366366366\\
-0.432432432432432	0.00252479881081081\\
-0.429429429429429	0.00252647395795796\\
-0.426426426426426	0.00252814910510511\\
-0.423423423423423	0.00252982425225225\\
-0.42042042042042	0.0025314993993994\\
-0.417417417417417	0.00253317454654655\\
-0.414414414414414	0.00253484969369369\\
-0.411411411411411	0.00253652484084084\\
-0.408408408408408	0.00253819998798799\\
-0.405405405405405	0.00253987513513514\\
-0.402402402402402	0.00254155028228228\\
-0.399399399399399	0.00254322542942943\\
-0.396396396396396	0.00254490057657658\\
-0.393393393393393	0.00254657572372372\\
-0.39039039039039	0.00254825087087087\\
-0.387387387387387	0.00254992601801802\\
-0.384384384384384	0.00255160116516517\\
-0.381381381381381	0.00255327631231231\\
-0.378378378378378	0.00255495145945946\\
-0.375375375375375	0.00255662660660661\\
-0.372372372372372	0.00255830175375375\\
-0.369369369369369	0.0025599769009009\\
-0.366366366366366	0.00256165204804805\\
-0.363363363363363	0.0025633271951952\\
-0.36036036036036	0.00256500234234234\\
-0.357357357357357	0.00256667748948949\\
-0.354354354354354	0.00256835263663664\\
-0.351351351351351	0.00257002778378378\\
-0.348348348348348	0.00257170293093093\\
-0.345345345345345	0.00257337807807808\\
-0.342342342342342	0.00257505322522523\\
-0.339339339339339	0.00257672837237237\\
-0.336336336336336	0.00257840351951952\\
-0.333333333333333	0.00258007866666667\\
-0.33033033033033	0.00258175381381381\\
-0.327327327327327	0.00258342896096096\\
-0.324324324324324	0.00258510410810811\\
-0.321321321321321	0.00258677925525526\\
-0.318318318318318	0.0025884544024024\\
-0.315315315315315	0.00259012954954955\\
-0.312312312312312	0.0025918046966967\\
-0.309309309309309	0.00259347984384384\\
-0.306306306306306	0.00259515499099099\\
-0.303303303303303	0.00259683013813814\\
-0.3003003003003	0.00259850528528529\\
-0.297297297297297	0.00260018043243243\\
-0.294294294294294	0.00260185557957958\\
-0.291291291291291	0.00260353072672673\\
-0.288288288288288	0.00260520587387387\\
-0.285285285285285	0.00260688102102102\\
-0.282282282282282	0.00260855616816817\\
-0.279279279279279	0.00261023131531532\\
-0.276276276276276	0.00261190646246246\\
-0.273273273273273	0.00261358160960961\\
-0.27027027027027	0.00261525675675676\\
-0.267267267267267	0.0026169319039039\\
-0.264264264264264	0.00261860705105105\\
-0.261261261261261	0.0026202821981982\\
-0.258258258258258	0.00262195734534535\\
-0.255255255255255	0.00262363249249249\\
-0.252252252252252	0.00262530763963964\\
-0.249249249249249	0.00262698278678679\\
-0.246246246246246	0.00262865793393393\\
-0.243243243243243	0.00263033308108108\\
-0.24024024024024	0.00263200822822823\\
-0.237237237237237	0.00263368337537538\\
-0.234234234234234	0.00263535852252252\\
-0.231231231231231	0.00263703366966967\\
-0.228228228228228	0.00263870881681682\\
-0.225225225225225	0.00264038396396396\\
-0.222222222222222	0.00264205911111111\\
-0.219219219219219	0.00264373425825826\\
-0.216216216216216	0.00264540940540541\\
-0.213213213213213	0.00264708455255255\\
-0.21021021021021	0.0026487596996997\\
-0.207207207207207	0.00265043484684685\\
-0.204204204204204	0.00265210999399399\\
-0.201201201201201	0.00265378514114114\\
-0.198198198198198	0.00265546028828829\\
-0.195195195195195	0.00265713543543544\\
-0.192192192192192	0.00265881058258258\\
-0.189189189189189	0.00266048572972973\\
-0.186186186186186	0.00266216087687688\\
-0.183183183183183	0.00266383602402402\\
-0.18018018018018	0.00266551117117117\\
-0.177177177177177	0.00266718631831832\\
-0.174174174174174	0.00266886146546547\\
-0.171171171171171	0.00267053661261261\\
-0.168168168168168	0.00267221175975976\\
-0.165165165165165	0.00267388690690691\\
-0.162162162162162	0.00267556205405405\\
-0.159159159159159	0.0026772372012012\\
-0.156156156156156	0.00267891234834835\\
-0.153153153153153	0.0026805874954955\\
-0.15015015015015	0.00268226264264264\\
-0.147147147147147	0.00268393778978979\\
-0.144144144144144	0.00268561293693694\\
-0.141141141141141	0.00268728808408408\\
-0.138138138138138	0.00268896323123123\\
-0.135135135135135	0.00269063837837838\\
-0.132132132132132	0.00269231352552553\\
-0.129129129129129	0.00269398867267267\\
-0.126126126126126	0.00269566381981982\\
-0.123123123123123	0.00269733896696697\\
-0.12012012012012	0.00269901411411411\\
-0.117117117117117	0.00270068926126126\\
-0.114114114114114	0.00270236440840841\\
-0.111111111111111	0.00270403955555556\\
-0.108108108108108	0.0027057147027027\\
-0.105105105105105	0.00270738984984985\\
-0.102102102102102	0.002709064996997\\
-0.0990990990990991	0.00271074014414414\\
-0.0960960960960962	0.00271241529129129\\
-0.0930930930930931	0.00271409043843844\\
-0.0900900900900901	0.00271576558558559\\
-0.0870870870870871	0.00271744073273273\\
-0.0840840840840841	0.00271911587987988\\
-0.081081081081081	0.00272079102702703\\
-0.0780780780780781	0.00272246617417417\\
-0.075075075075075	0.00272414132132132\\
-0.0720720720720721	0.00272581646846847\\
-0.0690690690690691	0.00272749161561562\\
-0.066066066066066	0.00272916676276276\\
-0.0630630630630631	0.00273084190990991\\
-0.06006006006006	0.00273251705705706\\
-0.0570570570570571	0.0027341922042042\\
-0.0540540540540541	0.00273586735135135\\
-0.051051051051051	0.0027375424984985\\
-0.0480480480480481	0.00273921764564565\\
-0.045045045045045	0.00274089279279279\\
-0.0420420420420421	0.00274256793993994\\
-0.039039039039039	0.00274424308708709\\
-0.036036036036036	0.00274591823423423\\
-0.0330330330330331	0.00274759338138138\\
-0.03003003003003	0.00274926852852853\\
-0.027027027027027	0.00275094367567568\\
-0.024024024024024	0.00275261882282282\\
-0.021021021021021	0.00275429396996997\\
-0.0180180180180181	0.00275596911711712\\
-0.015015015015015	0.00275764426426426\\
-0.012012012012012	0.00275931941141141\\
-0.00900900900900903	0.00276099455855856\\
-0.00600600600600598	0.00276266970570571\\
-0.00300300300300305	0.00276434485285285\\
0	0.00276602\\
0.00300300300300305	0.00276769514714715\\
0.00600600600600609	0.00276937029429429\\
0.00900900900900892	0.00277104544144144\\
0.012012012012012	0.00277272058858859\\
0.015015015015015	0.00277439573573574\\
0.0180180180180181	0.00277607088288288\\
0.0210210210210211	0.00277774603003003\\
0.0240240240240239	0.00277942117717718\\
0.027027027027027	0.00278109632432432\\
0.03003003003003	0.00278277147147147\\
0.0330330330330331	0.00278444661861862\\
0.0360360360360361	0.00278612176576577\\
0.0390390390390389	0.00278779691291291\\
0.042042042042042	0.00278947206006006\\
0.045045045045045	0.00279114720720721\\
0.0480480480480481	0.00279282235435435\\
0.0510510510510511	0.0027944975015015\\
0.0540540540540539	0.00279617264864865\\
0.057057057057057	0.0027978477957958\\
0.06006006006006	0.00279952294294294\\
0.0630630630630631	0.00280119809009009\\
0.0660660660660661	0.00280287323723724\\
0.0690690690690692	0.00280454838438438\\
0.072072072072072	0.00280622353153153\\
0.075075075075075	0.00280789867867868\\
0.0780780780780781	0.00280957382582583\\
0.0810810810810811	0.00281124897297297\\
0.0840840840840842	0.00281292412012012\\
0.087087087087087	0.00281459926726727\\
0.0900900900900901	0.00281627441441441\\
0.0930930930930931	0.00281794956156156\\
0.0960960960960962	0.00281962470870871\\
0.0990990990990992	0.00282129985585586\\
0.102102102102102	0.002822975003003\\
0.105105105105105	0.00282465015015015\\
0.108108108108108	0.0028263252972973\\
0.111111111111111	0.00282800044444444\\
0.114114114114114	0.00282967559159159\\
0.117117117117117	0.00283135073873874\\
0.12012012012012	0.00283302588588589\\
0.123123123123123	0.00283470103303303\\
0.126126126126126	0.00283637618018018\\
0.129129129129129	0.00283805132732733\\
0.132132132132132	0.00283972647447447\\
0.135135135135135	0.00284140162162162\\
0.138138138138138	0.00284307676876877\\
0.141141141141141	0.00284475191591592\\
0.144144144144144	0.00284642706306306\\
0.147147147147147	0.00284810221021021\\
0.15015015015015	0.00284977735735736\\
0.153153153153153	0.0028514525045045\\
0.156156156156156	0.00285312765165165\\
0.159159159159159	0.0028548027987988\\
0.162162162162162	0.00285647794594595\\
0.165165165165165	0.00285815309309309\\
0.168168168168168	0.00285982824024024\\
0.171171171171171	0.00286150338738739\\
0.174174174174174	0.00286317853453453\\
0.177177177177177	0.00286485368168168\\
0.18018018018018	0.00286652882882883\\
0.183183183183183	0.00286820397597598\\
0.186186186186186	0.00286987912312312\\
0.189189189189189	0.00287155427027027\\
0.192192192192192	0.00287322941741742\\
0.195195195195195	0.00287490456456456\\
0.198198198198198	0.00287657971171171\\
0.201201201201201	0.00287825485885886\\
0.204204204204204	0.00287993000600601\\
0.207207207207207	0.00288160515315315\\
0.21021021021021	0.0028832803003003\\
0.213213213213213	0.00288495544744745\\
0.216216216216216	0.00288663059459459\\
0.219219219219219	0.00288830574174174\\
0.222222222222222	0.00288998088888889\\
0.225225225225225	0.00289165603603604\\
0.228228228228228	0.00289333118318318\\
0.231231231231231	0.00289500633033033\\
0.234234234234234	0.00289668147747748\\
0.237237237237237	0.00289835662462462\\
0.24024024024024	0.00290003177177177\\
0.243243243243243	0.00290170691891892\\
0.246246246246246	0.00290338206606607\\
0.249249249249249	0.00290505721321321\\
0.252252252252252	0.00290673236036036\\
0.255255255255255	0.00290840750750751\\
0.258258258258258	0.00291008265465465\\
0.261261261261261	0.0029117578018018\\
0.264264264264264	0.00291343294894895\\
0.267267267267267	0.0029151080960961\\
0.27027027027027	0.00291678324324324\\
0.273273273273273	0.00291845839039039\\
0.276276276276276	0.00292013353753754\\
0.279279279279279	0.00292180868468468\\
0.282282282282282	0.00292348383183183\\
0.285285285285285	0.00292515897897898\\
0.288288288288288	0.00292683412612613\\
0.291291291291291	0.00292850927327327\\
0.294294294294294	0.00293018442042042\\
0.297297297297297	0.00293185956756757\\
0.3003003003003	0.00293353471471471\\
0.303303303303303	0.00293520986186186\\
0.306306306306306	0.00293688500900901\\
0.309309309309309	0.00293856015615616\\
0.312312312312312	0.0029402353033033\\
0.315315315315315	0.00294191045045045\\
0.318318318318318	0.0029435855975976\\
0.321321321321321	0.00294526074474474\\
0.324324324324324	0.00294693589189189\\
0.327327327327327	0.00294861103903904\\
0.33033033033033	0.00295028618618619\\
0.333333333333333	0.00295196133333333\\
0.336336336336336	0.00295363648048048\\
0.339339339339339	0.00295531162762763\\
0.342342342342342	0.00295698677477477\\
0.345345345345345	0.00295866192192192\\
0.348348348348348	0.00296033706906907\\
0.351351351351351	0.00296201221621622\\
0.354354354354354	0.00296368736336336\\
0.357357357357357	0.00296536251051051\\
0.36036036036036	0.00296703765765766\\
0.363363363363363	0.0029687128048048\\
0.366366366366366	0.00297038795195195\\
0.369369369369369	0.0029720630990991\\
0.372372372372372	0.00297373824624625\\
0.375375375375375	0.00297541339339339\\
0.378378378378378	0.00297708854054054\\
0.381381381381381	0.00297876368768769\\
0.384384384384384	0.00298043883483484\\
0.387387387387387	0.00298211398198198\\
0.39039039039039	0.00298378912912913\\
0.393393393393393	0.00298546427627628\\
0.396396396396396	0.00298713942342342\\
0.399399399399399	0.00298881457057057\\
0.402402402402402	0.00299048971771772\\
0.405405405405405	0.00299216486486486\\
0.408408408408408	0.00299384001201201\\
0.411411411411411	0.00299551515915916\\
0.414414414414414	0.00299719030630631\\
0.417417417417417	0.00299886545345345\\
0.42042042042042	0.0030005406006006\\
0.423423423423423	0.00300221574774775\\
0.426426426426426	0.0030038908948949\\
0.429429429429429	0.00300556604204204\\
0.432432432432432	0.00300724118918919\\
0.435435435435436	0.00300891633633634\\
0.438438438438439	0.00301059148348348\\
0.441441441441441	0.00301226663063063\\
0.444444444444444	0.00301394177777778\\
0.447447447447447	0.00301561692492492\\
0.450450450450451	0.00301729207207207\\
0.453453453453454	0.00301896721921922\\
0.456456456456456	0.00302064236636637\\
0.459459459459459	0.00302231751351351\\
0.462462462462462	0.00302399266066066\\
0.465465465465466	0.00302566780780781\\
0.468468468468469	0.00302734295495496\\
0.471471471471471	0.0030290181021021\\
0.474474474474474	0.00303069324924925\\
0.477477477477477	0.0030323683963964\\
0.480480480480481	0.00303404354354354\\
0.483483483483484	0.00303571869069069\\
0.486486486486486	0.00303739383783784\\
0.489489489489489	0.00303906898498498\\
0.492492492492492	0.00304074413213213\\
0.495495495495496	0.00304241927927928\\
0.498498498498499	0.00304409442642643\\
0.501501501501501	0.00304576957357357\\
0.504504504504504	0.00304744472072072\\
0.507507507507508	0.00304911986786787\\
0.510510510510511	0.00305079501501502\\
0.513513513513514	0.00305247016216216\\
0.516516516516516	0.00305414530930931\\
0.519519519519519	0.00305582045645646\\
0.522522522522523	0.0030574956036036\\
0.525525525525526	0.00305917075075075\\
0.528528528528529	0.0030608458978979\\
0.531531531531531	0.00306252104504505\\
0.534534534534534	0.00306419619219219\\
0.537537537537538	0.00306587133933934\\
0.540540540540541	0.00306754648648649\\
0.543543543543544	0.00306922163363363\\
0.546546546546546	0.00307089678078078\\
0.549549549549549	0.00307257192792793\\
0.552552552552553	0.00307424707507508\\
0.555555555555556	0.00307592222222222\\
0.558558558558559	0.00307759736936937\\
0.561561561561561	0.00307927251651652\\
0.564564564564564	0.00308094766366366\\
0.567567567567568	0.00308262281081081\\
0.570570570570571	0.00308429795795796\\
0.573573573573574	0.00308597310510511\\
0.576576576576577	0.00308764825225225\\
0.57957957957958	0.0030893233993994\\
0.582582582582583	0.00309099854654655\\
0.585585585585586	0.00309267369369369\\
0.588588588588589	0.00309434884084084\\
0.591591591591592	0.00309602398798799\\
0.594594594594595	0.00309769913513514\\
0.597597597597598	0.00309937428228228\\
0.600600600600601	0.00310104942942943\\
0.603603603603604	0.00310272457657658\\
0.606606606606607	0.00310439972372372\\
0.60960960960961	0.00310607487087087\\
0.612612612612613	0.00310775001801802\\
0.615615615615616	0.00310942516516517\\
0.618618618618619	0.00311110031231231\\
0.621621621621622	0.00311277545945946\\
0.624624624624625	0.00311445060660661\\
0.627627627627628	0.00311612575375375\\
0.630630630630631	0.0031178009009009\\
0.633633633633634	0.00311947604804805\\
0.636636636636637	0.0031211511951952\\
0.63963963963964	0.00312282634234234\\
0.642642642642643	0.00312450148948949\\
0.645645645645646	0.00312617663663664\\
0.648648648648649	0.00312785178378378\\
0.651651651651652	0.00312952693093093\\
0.654654654654655	0.00313120207807808\\
0.657657657657658	0.00313287722522523\\
0.660660660660661	0.00313455237237237\\
0.663663663663664	0.00313622751951952\\
0.666666666666667	0.00313790266666667\\
0.66966966966967	0.00313957781381381\\
0.672672672672673	0.00314125296096096\\
0.675675675675676	0.00314292810810811\\
0.678678678678679	0.00314460325525526\\
0.681681681681682	0.0031462784024024\\
0.684684684684685	0.00314795354954955\\
0.687687687687688	0.0031496286966967\\
0.690690690690691	0.00315130384384384\\
0.693693693693694	0.00315297899099099\\
0.696696696696697	0.00315465413813814\\
0.6996996996997	0.00315632928528529\\
0.702702702702703	0.00315800443243243\\
0.705705705705706	0.00315967957957958\\
0.708708708708709	0.00316135472672673\\
0.711711711711712	0.00316302987387387\\
0.714714714714715	0.00316470502102102\\
0.717717717717718	0.00316638016816817\\
0.720720720720721	0.00316805531531532\\
0.723723723723724	0.00316973046246246\\
0.726726726726727	0.00317140560960961\\
0.72972972972973	0.00317308075675676\\
0.732732732732733	0.0031747559039039\\
0.735735735735736	0.00317643105105105\\
0.738738738738739	0.0031781061981982\\
0.741741741741742	0.00317978134534535\\
0.744744744744745	0.00318145649249249\\
0.747747747747748	0.00318313163963964\\
0.750750750750751	0.00318480678678679\\
0.753753753753754	0.00318648193393393\\
0.756756756756757	0.00318815708108108\\
0.75975975975976	0.00318983222822823\\
0.762762762762763	0.00319150737537538\\
0.765765765765766	0.00319318252252252\\
0.768768768768769	0.00319485766966967\\
0.771771771771772	0.00319653281681682\\
0.774774774774775	0.00319820796396396\\
0.777777777777778	0.00319988311111111\\
0.780780780780781	0.00320155825825826\\
0.783783783783784	0.00320323340540541\\
0.786786786786787	0.00320490855255255\\
0.78978978978979	0.0032065836996997\\
0.792792792792793	0.00320825884684685\\
0.795795795795796	0.00320993399399399\\
0.798798798798799	0.00321160914114114\\
0.801801801801802	0.00321328428828829\\
0.804804804804805	0.00321495943543544\\
0.807807807807808	0.00321663458258258\\
0.810810810810811	0.00321830972972973\\
0.813813813813814	0.00321998487687688\\
0.816816816816817	0.00322166002402402\\
0.81981981981982	0.00322333517117117\\
0.822822822822823	0.00322501031831832\\
0.825825825825826	0.00322668546546547\\
0.828828828828829	0.00322836061261261\\
0.831831831831832	0.00323003575975976\\
0.834834834834835	0.00323171090690691\\
0.837837837837838	0.00323338605405405\\
0.840840840840841	0.0032350612012012\\
0.843843843843844	0.00323673634834835\\
0.846846846846847	0.0032384114954955\\
0.84984984984985	0.00324008664264264\\
0.852852852852853	0.00324176178978979\\
0.855855855855856	0.00324343693693694\\
0.858858858858859	0.00324511208408408\\
0.861861861861862	0.00324678723123123\\
0.864864864864865	0.00324846237837838\\
0.867867867867868	0.00325013752552553\\
0.870870870870871	0.00325181267267267\\
0.873873873873874	0.00325348781981982\\
0.876876876876877	0.00325516296696697\\
0.87987987987988	0.00325683811411411\\
0.882882882882883	0.00325851326126126\\
0.885885885885886	0.00326018840840841\\
0.888888888888889	0.00326186355555556\\
0.891891891891892	0.0032635387027027\\
0.894894894894895	0.00326521384984985\\
0.897897897897898	0.003266888996997\\
0.900900900900901	0.00326856414414414\\
0.903903903903904	0.00327023929129129\\
0.906906906906907	0.00327191443843844\\
0.90990990990991	0.00327358958558559\\
0.912912912912913	0.00327526473273273\\
0.915915915915916	0.00327693987987988\\
0.918918918918919	0.00327861502702703\\
0.921921921921922	0.00328029017417417\\
0.924924924924925	0.00328196532132132\\
0.927927927927928	0.00328364046846847\\
0.930930930930931	0.00328531561561562\\
0.933933933933934	0.00328699076276276\\
0.936936936936937	0.00328866590990991\\
0.93993993993994	0.00329034105705706\\
0.942942942942943	0.0032920162042042\\
0.945945945945946	0.00329369135135135\\
0.948948948948949	0.0032953664984985\\
0.951951951951952	0.00329704164564565\\
0.954954954954955	0.00329871679279279\\
0.957957957957958	0.00330039193993994\\
0.960960960960961	0.00330206708708709\\
0.963963963963964	0.00330374223423423\\
0.966966966966967	0.00330541738138138\\
0.96996996996997	0.00330709252852853\\
0.972972972972973	0.00330876767567568\\
0.975975975975976	0.00331044282282282\\
0.978978978978979	0.00331211796996997\\
0.981981981981982	0.00331379311711712\\
0.984984984984985	0.00331546826426426\\
0.987987987987988	0.00331714341141141\\
0.990990990990991	0.00331881855855856\\
0.993993993993994	0.00332049370570571\\
0.996996996996997	0.00332216885285285\\
1	0.003323844\\
1.003003003003	0.00332551914714715\\
1.00600600600601	0.00332719429429429\\
1.00900900900901	0.00332886944144144\\
1.01201201201201	0.00333054458858859\\
1.01501501501502	0.00333221973573574\\
1.01801801801802	0.00333389488288288\\
1.02102102102102	0.00333557003003003\\
1.02402402402402	0.00333724517717718\\
1.02702702702703	0.00333892032432432\\
1.03003003003003	0.00334059547147147\\
1.03303303303303	0.00334227061861862\\
1.03603603603604	0.00334394576576577\\
1.03903903903904	0.00334562091291291\\
1.04204204204204	0.00334729606006006\\
1.04504504504505	0.00334897120720721\\
1.04804804804805	0.00335064635435435\\
1.05105105105105	0.0033523215015015\\
1.05405405405405	0.00335399664864865\\
1.05705705705706	0.0033556717957958\\
1.06006006006006	0.00335734694294294\\
1.06306306306306	0.00335902209009009\\
1.06606606606607	0.00336069723723724\\
1.06906906906907	0.00336237238438438\\
1.07207207207207	0.00336404753153153\\
1.07507507507508	0.00336572267867868\\
1.07807807807808	0.00336739782582583\\
1.08108108108108	0.00336907297297297\\
1.08408408408408	0.00337074812012012\\
1.08708708708709	0.00337242326726727\\
1.09009009009009	0.00337409841441441\\
1.09309309309309	0.00337577356156156\\
1.0960960960961	0.00337744870870871\\
1.0990990990991	0.00337912385585586\\
1.1021021021021	0.003380799003003\\
1.10510510510511	0.00338247415015015\\
1.10810810810811	0.0033841492972973\\
1.11111111111111	0.00338582444444444\\
1.11411411411411	0.00338749959159159\\
1.11711711711712	0.00338917473873874\\
1.12012012012012	0.00339084988588589\\
1.12312312312312	0.00339252503303303\\
1.12612612612613	0.00339420018018018\\
1.12912912912913	0.00339587532732733\\
1.13213213213213	0.00339755047447447\\
1.13513513513514	0.00339922562162162\\
1.13813813813814	0.00340090076876877\\
1.14114114114114	0.00340257591591592\\
1.14414414414414	0.00340425106306306\\
1.14714714714715	0.00340592621021021\\
1.15015015015015	0.00340760135735736\\
1.15315315315315	0.0034092765045045\\
1.15615615615616	0.00341095165165165\\
1.15915915915916	0.0034126267987988\\
1.16216216216216	0.00341430194594595\\
1.16516516516517	0.00341597709309309\\
1.16816816816817	0.00341765224024024\\
1.17117117117117	0.00341932738738739\\
1.17417417417417	0.00342100253453453\\
1.17717717717718	0.00342267768168168\\
1.18018018018018	0.00342435282882883\\
1.18318318318318	0.00342602797597598\\
1.18618618618619	0.00342770312312312\\
1.18918918918919	0.00342937827027027\\
1.19219219219219	0.00343105341741742\\
1.1951951951952	0.00343272856456456\\
1.1981981981982	0.00343440371171171\\
1.2012012012012	0.00343607885885886\\
1.2042042042042	0.00343775400600601\\
1.20720720720721	0.00343942915315315\\
1.21021021021021	0.0034411043003003\\
1.21321321321321	0.00344277944744745\\
1.21621621621622	0.00344445459459459\\
1.21921921921922	0.00344612974174174\\
1.22222222222222	0.00344780488888889\\
1.22522522522523	0.00344948003603604\\
1.22822822822823	0.00345115518318318\\
1.23123123123123	0.00345283033033033\\
1.23423423423423	0.00345450547747748\\
1.23723723723724	0.00345618062462462\\
1.24024024024024	0.00345785577177177\\
1.24324324324324	0.00345953091891892\\
1.24624624624625	0.00346120606606607\\
1.24924924924925	0.00346288121321321\\
1.25225225225225	0.00346455636036036\\
1.25525525525526	0.00346623150750751\\
1.25825825825826	0.00346790665465465\\
1.26126126126126	0.0034695818018018\\
1.26426426426426	0.00347125694894895\\
1.26726726726727	0.0034729320960961\\
1.27027027027027	0.00347460724324324\\
1.27327327327327	0.00347628239039039\\
1.27627627627628	0.00347795753753754\\
1.27927927927928	0.00347963268468468\\
1.28228228228228	0.00348130783183183\\
1.28528528528529	0.00348298297897898\\
1.28828828828829	0.00348465812612613\\
1.29129129129129	0.00348633327327327\\
1.29429429429429	0.00348800842042042\\
1.2972972972973	0.00348968356756757\\
1.3003003003003	0.00349135871471471\\
1.3033033033033	0.00349303386186186\\
1.30630630630631	0.00349470900900901\\
1.30930930930931	0.00349638415615616\\
1.31231231231231	0.0034980593033033\\
1.31531531531532	0.00349973445045045\\
1.31831831831832	0.0035014095975976\\
1.32132132132132	0.00350308474474474\\
1.32432432432432	0.00350475989189189\\
1.32732732732733	0.00350643503903904\\
1.33033033033033	0.00350811018618619\\
1.33333333333333	0.00350978533333333\\
1.33633633633634	0.00351146048048048\\
1.33933933933934	0.00351313562762763\\
1.34234234234234	0.00351481077477477\\
1.34534534534535	0.00351648592192192\\
1.34834834834835	0.00351816106906907\\
1.35135135135135	0.00351983621621622\\
1.35435435435435	0.00352151136336336\\
1.35735735735736	0.00352318651051051\\
1.36036036036036	0.00352486165765766\\
1.36336336336336	0.00352653680480481\\
1.36636636636637	0.00352821195195195\\
1.36936936936937	0.0035298870990991\\
1.37237237237237	0.00353156224624625\\
1.37537537537538	0.00353323739339339\\
1.37837837837838	0.00353491254054054\\
1.38138138138138	0.00353658768768769\\
1.38438438438438	0.00353826283483483\\
1.38738738738739	0.00353993798198198\\
1.39039039039039	0.00354161312912913\\
1.39339339339339	0.00354328827627628\\
1.3963963963964	0.00354496342342342\\
1.3993993993994	0.00354663857057057\\
1.4024024024024	0.00354831371771772\\
1.40540540540541	0.00354998886486487\\
1.40840840840841	0.00355166401201201\\
1.41141141141141	0.00355333915915916\\
1.41441441441441	0.00355501430630631\\
1.41741741741742	0.00355668945345345\\
1.42042042042042	0.0035583646006006\\
1.42342342342342	0.00356003974774775\\
1.42642642642643	0.00356171489489489\\
1.42942942942943	0.00356339004204204\\
1.43243243243243	0.00356506518918919\\
1.43543543543544	0.00356674033633634\\
1.43843843843844	0.00356841548348348\\
1.44144144144144	0.00357009063063063\\
1.44444444444444	0.00357176577777778\\
1.44744744744745	0.00357344092492493\\
1.45045045045045	0.00357511607207207\\
1.45345345345345	0.00357679121921922\\
1.45645645645646	0.00357846636636637\\
1.45945945945946	0.00358014151351351\\
1.46246246246246	0.00358181666066066\\
1.46546546546547	0.00358349180780781\\
1.46846846846847	0.00358516695495496\\
1.47147147147147	0.0035868421021021\\
1.47447447447447	0.00358851724924925\\
1.47747747747748	0.0035901923963964\\
1.48048048048048	0.00359186754354354\\
1.48348348348348	0.00359354269069069\\
1.48648648648649	0.00359521783783784\\
1.48948948948949	0.00359689298498499\\
1.49249249249249	0.00359856813213213\\
1.4954954954955	0.00360024327927928\\
1.4984984984985	0.00360191842642643\\
1.5015015015015	0.00360359357357357\\
1.5045045045045	0.00360526872072072\\
1.50750750750751	0.00360694386786787\\
1.51051051051051	0.00360861901501501\\
1.51351351351351	0.00361029416216216\\
1.51651651651652	0.00361196930930931\\
1.51951951951952	0.00361364445645646\\
1.52252252252252	0.0036153196036036\\
1.52552552552553	0.00361699475075075\\
1.52852852852853	0.0036186698978979\\
1.53153153153153	0.00362034504504505\\
1.53453453453453	0.00362202019219219\\
1.53753753753754	0.00362369533933934\\
1.54054054054054	0.00362537048648649\\
1.54354354354354	0.00362704563363363\\
1.54654654654655	0.00362872078078078\\
1.54954954954955	0.00363039592792793\\
1.55255255255255	0.00363207107507508\\
1.55555555555556	0.00363374622222222\\
1.55855855855856	0.00363542136936937\\
1.56156156156156	0.00363709651651652\\
1.56456456456456	0.00363877166366366\\
1.56756756756757	0.00364044681081081\\
1.57057057057057	0.00364212195795796\\
1.57357357357357	0.00364379710510511\\
1.57657657657658	0.00364547225225225\\
1.57957957957958	0.0036471473993994\\
1.58258258258258	0.00364882254654655\\
1.58558558558559	0.00365049769369369\\
1.58858858858859	0.00365217284084084\\
1.59159159159159	0.00365384798798799\\
1.59459459459459	0.00365552313513514\\
1.5975975975976	0.00365719828228228\\
1.6006006006006	0.00365887342942943\\
1.6036036036036	0.00366054857657658\\
1.60660660660661	0.00366222372372372\\
1.60960960960961	0.00366389887087087\\
1.61261261261261	0.00366557401801802\\
1.61561561561562	0.00366724916516516\\
1.61861861861862	0.00366892431231231\\
1.62162162162162	0.00367059945945946\\
1.62462462462462	0.00367227460660661\\
1.62762762762763	0.00367394975375375\\
1.63063063063063	0.0036756249009009\\
1.63363363363363	0.00367730004804805\\
1.63663663663664	0.0036789751951952\\
1.63963963963964	0.00368065034234234\\
1.64264264264264	0.00368232548948949\\
1.64564564564565	0.00368400063663664\\
1.64864864864865	0.00368567578378378\\
1.65165165165165	0.00368735093093093\\
1.65465465465465	0.00368902607807808\\
1.65765765765766	0.00369070122522522\\
1.66066066066066	0.00369237637237237\\
1.66366366366366	0.00369405151951952\\
1.66666666666667	0.00369572666666667\\
1.66966966966967	0.00369740181381381\\
1.67267267267267	0.00369907696096096\\
1.67567567567568	0.00370075210810811\\
1.67867867867868	0.00370242725525526\\
1.68168168168168	0.0037041024024024\\
1.68468468468468	0.00370577754954955\\
1.68768768768769	0.0037074526966967\\
1.69069069069069	0.00370912784384384\\
1.69369369369369	0.00371080299099099\\
1.6966966966967	0.00371247813813814\\
1.6996996996997	0.00371415328528529\\
1.7027027027027	0.00371582843243243\\
1.70570570570571	0.00371750357957958\\
1.70870870870871	0.00371917872672673\\
1.71171171171171	0.00372085387387387\\
1.71471471471471	0.00372252902102102\\
1.71771771771772	0.00372420416816817\\
1.72072072072072	0.00372587931531532\\
1.72372372372372	0.00372755446246246\\
1.72672672672673	0.00372922960960961\\
1.72972972972973	0.00373090475675676\\
1.73273273273273	0.0037325799039039\\
1.73573573573574	0.00373425505105105\\
1.73873873873874	0.0037359301981982\\
1.74174174174174	0.00373760534534535\\
1.74474474474474	0.00373928049249249\\
1.74774774774775	0.00374095563963964\\
1.75075075075075	0.00374263078678679\\
1.75375375375375	0.00374430593393393\\
1.75675675675676	0.00374598108108108\\
1.75975975975976	0.00374765622822823\\
1.76276276276276	0.00374933137537538\\
1.76576576576577	0.00375100652252252\\
1.76876876876877	0.00375268166966967\\
1.77177177177177	0.00375435681681682\\
1.77477477477477	0.00375603196396396\\
1.77777777777778	0.00375770711111111\\
1.78078078078078	0.00375938225825826\\
1.78378378378378	0.00376105740540541\\
1.78678678678679	0.00376273255255255\\
1.78978978978979	0.0037644076996997\\
1.79279279279279	0.00376608284684685\\
1.7957957957958	0.00376775799399399\\
1.7987987987988	0.00376943314114114\\
1.8018018018018	0.00377110828828829\\
1.8048048048048	0.00377278343543544\\
1.80780780780781	0.00377445858258258\\
1.81081081081081	0.00377613372972973\\
1.81381381381381	0.00377780887687688\\
1.81681681681682	0.00377948402402402\\
1.81981981981982	0.00378115917117117\\
1.82282282282282	0.00378283431831832\\
1.82582582582583	0.00378450946546547\\
1.82882882882883	0.00378618461261261\\
1.83183183183183	0.00378785975975976\\
1.83483483483483	0.00378953490690691\\
1.83783783783784	0.00379121005405405\\
1.84084084084084	0.0037928852012012\\
1.84384384384384	0.00379456034834835\\
1.84684684684685	0.0037962354954955\\
1.84984984984985	0.00379791064264264\\
1.85285285285285	0.00379958578978979\\
1.85585585585586	0.00380126093693694\\
1.85885885885886	0.00380293608408408\\
1.86186186186186	0.00380461123123123\\
1.86486486486486	0.00380628637837838\\
1.86786786786787	0.00380796152552553\\
1.87087087087087	0.00380963667267267\\
1.87387387387387	0.00381131181981982\\
1.87687687687688	0.00381298696696697\\
1.87987987987988	0.00381466211411411\\
1.88288288288288	0.00381633726126126\\
1.88588588588589	0.00381801240840841\\
1.88888888888889	0.00381968755555556\\
1.89189189189189	0.0038213627027027\\
1.89489489489489	0.00382303784984985\\
1.8978978978979	0.003824712996997\\
1.9009009009009	0.00382638814414414\\
1.9039039039039	0.00382806329129129\\
1.90690690690691	0.00382973843843844\\
1.90990990990991	0.00383141358558559\\
1.91291291291291	0.00383308873273273\\
1.91591591591592	0.00383476387987988\\
1.91891891891892	0.00383643902702703\\
1.92192192192192	0.00383811417417417\\
1.92492492492492	0.00383978932132132\\
1.92792792792793	0.00384146446846847\\
1.93093093093093	0.00384313961561562\\
1.93393393393393	0.00384481476276276\\
1.93693693693694	0.00384648990990991\\
1.93993993993994	0.00384816505705706\\
1.94294294294294	0.0038498402042042\\
1.94594594594595	0.00385151535135135\\
1.94894894894895	0.0038531904984985\\
1.95195195195195	0.00385486564564565\\
1.95495495495495	0.00385654079279279\\
1.95795795795796	0.00385821593993994\\
1.96096096096096	0.00385989108708709\\
1.96396396396396	0.00386156623423423\\
1.96696696696697	0.00386324138138138\\
1.96996996996997	0.00386491652852853\\
1.97297297297297	0.00386659167567568\\
1.97597597597598	0.00386826682282282\\
1.97897897897898	0.00386994196996997\\
1.98198198198198	0.00387161711711712\\
1.98498498498498	0.00387329226426426\\
1.98798798798799	0.00387496741141141\\
1.99099099099099	0.00387664255855856\\
1.99399399399399	0.00387831770570571\\
1.996996996997	0.00387999285285285\\
2	0.003881668\\
};
\addlegendentry{Linear}

\end{axis}
\end{tikzpicture}%
\end{figure}
For the cooling curve we get the following regression parameters

\noindent For cubic fit:
\begin{align*}
    A &= 2.82265 \times 10^{-3} \\
    B &= 0.298851 \times 10^{-3} \\
    C &= 0.236043 \times 10^{-3} \\
    R^{2}&=0.9856
\end{align*}
For linear fit:
\begin{align*}
    a_0 &= 2.76602 \times 10^{-3} \\
    a_1 &= 0.557824 \times 10^{-3} \\
    R^{2}&=0.9422
\end{align*}
\section*{Result}
We have verified that the resistance of the thermistor depends on the temperature wrt the Steinhart-Hart Equation. The thermistor is linear for small temperature ranges.
\end{document}